\documentclass[12pt]{amsart}
\usepackage[T1]{fontenc}
\usepackage[utf8]{inputenc}
\usepackage{lmodern}
\usepackage[letterpaper,top=1.0in,bottom=1.0in,left=1.18in,right=1.18in]{geometry}
\usepackage{amsmath,amssymb,amsthm,mathtools}
\usepackage{microtype}

\setlength{\parindent}{15pt}
\setlength{\parskip}{0pt}
\setlength{\abovedisplayskip}{8pt plus 2pt minus 4pt}
\setlength{\belowdisplayskip}{8pt plus 2pt minus 4pt}
\setlength{\abovedisplayshortskip}{6pt plus 2pt minus 3pt}
\setlength{\belowdisplayshortskip}{6pt plus 2pt minus 3pt}
\DeclareMathOperator{\len}{len}
\linespread{1.01}
\title{A Short Path Joining Two Zeros Inside a Polynomial Lemniscate}


\theoremstyle{definition}
\newtheorem{theorem}{Theorem}
\newtheorem{lemma}[theorem]{Lemma}
\newtheorem{proposition}[theorem]{Proposition}
\newtheorem{corollary}[theorem]{Corollary}
\newtheorem*{componentlemma}{Component lemma}
\newtheorem*{auxlemma}{Auxiliary lemma (flux across a regular superlevel component)}
\theoremstyle{remark}
\newtheorem{remark}[theorem]{Remark}
\newtheorem*{remarkplain}{Remark}

\begin{document}

\begin{abstract}
Let
\[
 f(z)=\prod_{j=1}^n (z-z_j)
\]
be a monic polynomial whose zeros are listed with multiplicity and satisfy $|z_j|<1$ for all $j$.
Assuming the known Erd\H{o}s--Herzog--Piranian component lemma that some connected component of
\[
 \Lambda_f:=\{z\in\mathbb{C}: |f(z)|<1\}
\]
contains at least two zeros of $f$, we prove that there exist two roots of $f$ which can be joined by
 a rectifiable curve of Euclidean length $<2$ contained in $\Lambda_f$. The proof combines a sharp
 integral estimate on each component of a polynomial lemniscate, an explicit two-parameter perturbation
 that makes the critical-level structure above a fixed base level Morse and excellent on a fixed
 collar, and a spanning-tree construction in the resulting Reeb decomposition.
\end{abstract}

\maketitle

\section{Introduction}

Let
\[
 f(z)=\prod_{j=1}^n (z-z_j), \qquad |z_j|<1 \quad (1\le j\le n),
\]
where the zeros are listed with multiplicity. We consider the filled unit lemniscate
\[
 \Lambda_f:=\{z\in\mathbb{C}: |f(z)|<1\}.
\]
The question addressed here is the following.

\medskip
\noindent Must there always exist a rectifiable curve of Euclidean length less than $2$, contained in
$\Lambda_f$, whose endpoints are two roots of $f$?
\medskip

\noindent Our main result answers this in the affirmative.

\begin{theorem}\label{thm:main}
\emph{Let
\[
 f(z)=\prod_{j=1}^n (z-z_j), \qquad |z_j|<1 \quad (1\le j\le n),
\]
where the zeros are listed with multiplicity. Then there exist indices $i,j\in\{1,\dots,n\}$ and a rectifiable
curve $\gamma\subset \Lambda_f$ joining $z_i$ to $z_j$ such that
\[
 \operatorname{len}(\gamma)<2.
\]
In particular, if $f$ is squarefree and $n\ge 2$, then the endpoints may be chosen to be distinct zeros of $f$.}
\end{theorem}

\begin{remark}
If $f$ has a multiple zero, the conclusion is immediate: two coincident roots may be
joined by the constant curve. Thus the only nontrivial case is when $f$ is squarefree and $n\ge 2$.
In the proof below we therefore reduce immediately to that case.
\end{remark}

The main geometric input we use is the following component lemma.

\begin{componentlemma}
Some connected component of $\Lambda_f$ contains at least two zeros of $f$.
\end{componentlemma}

\medskip
\noindent
This is known due to work of Erd\H{o}s, Herzog, and Piranian; see \cite{EHP58} and the later references listed below. 
In addition, we use standard results from planar potential theory, planar topology, and Morse theory: Carleman's inequality, P\'olya's inequality, the Jordan--Schoenflies theorem, and the local Morse normal form for a nondegenerate saddle.
\noindent The proof of Theorem~\ref{thm:main} has three parts.
\begin{enumerate}
\item For any component $U\subset \Lambda_f$ containing $m$ zeros of $f$, we prove the sharp integral estimate
\[
 \int_U \frac{|f'(z)|}{|f(z)|}\,dA(z)\le 2\pi m.
\]
\item Above a fixed base level we construct, under a genericity hypothesis on the critical points, a finite
      tree spanning the relevant roots, with total length controlled by the tail integral of the level lengths.
\item A small explicit perturbation of the form $f\mapsto f+\lambda z+\beta$ makes the critical-point configuration
      generic on a fixed collar, while preserving the relevant component and a fixed positive slack below the threshold $2$.
\end{enumerate}

\section{An analytic estimate on a component}

Let $U$ be a connected component of $\Lambda_f$, and suppose that $U$ contains exactly $m\ge 1$ zeros of $f$,
counted with multiplicity. Set
\[
 u:=-\log|f| \quad\text{on } U.
\]
Then $u>0$ on $U$, $u=0$ on $\partial U$, and $u(z)\to +\infty$ at the zeros of $f$ in $U$. For $t\ge 0$ define
\[
 K_t:=\{z\in U:u(z)\ge t\}, \qquad A(t):=\operatorname{Area}(K_t).
\]
For regular $t$ write
\[
 P(t):=\mathcal H^1(\{z\in U:u(z)=t\}).
\]
We also set
\[
 I(U):=\int_U \frac{|f'(z)|}{|f(z)|}\,dA(z)=\int_U |\nabla u|\,dA.
\]

\medskip
\begin{lemma}\label{lem:level-set-measure-zero}
\emph{Let $p$ be a nonconstant polynomial and let $r>0$. Then the level set
\[
 \{z\in\mathbb{C}: |p(z)|=r\}
\]
has planar Lebesgue measure zero. Consequently, if $E$ is a connected component of
\[
 \{z\in\mathbb{C}: |p(z)|<r\},
\]
then $\partial E\subset\{|p|=r\}$ and
\[
 \operatorname{Area}(E)=\operatorname{Area}(\overline E).
\]}
\end{lemma}

\begin{proof} Set
\[
S:=\{z\in \mathbb{C}: |p(z)|=r,\ p'(z)\neq 0\}.
\]
The function $F(z):=|p(z)|^2-r^2$ is real-analytic on $\mathbb{C}$, and on $S$ its gradient is nonzero because,
after identifying $\mathbb{C}$ with $\mathbb{R}^2$,
\[
\nabla F(z)=2\,\overline{p(z)}\,p'(z).
\]
Therefore the implicit function theorem shows that $S$ is a one-dimensional real-analytic submanifold
of $\mathbb{C}$, hence it has planar measure zero. The complement
\[
\{|p|=r\}\setminus S
\]
is contained in the finite set $\{p'=0\}$, because $r>0$ excludes the zeros of $p$. Thus $\{|p|=r\}$ has
planar measure zero.

\indent If $E$ is a connected component of $\{|p|<r\}$, then $\partial E\subset\{|p|=r\}$. Hence
\[
\operatorname{Area}(\overline E\setminus E)=0,
\]
so $\operatorname{Area}(E)=\operatorname{Area}(\overline E)$. \end{proof}

\medskip
\begin{auxlemma}
Let $s>0$ be a regular value of $u=-\log|f|$, and let $G$ be a connected component of
\[
\{z\in U:u(z)>s\}.
\]
Suppose that the zeros of $f$ in $G$ are $a_1,\dots,a_r$, with multiplicities $m_1,\dots,m_r$, and set
\[
M:=m_1+\cdots+m_r.
\]
Then
\[
\int_{\partial G} |\nabla u|\,ds = 2\pi M.
\]
\end{auxlemma}

\begin{proof} Because $s$ is a regular value, $\partial G\subset\{u=s\}$ is a smooth compact $1$-manifold. For $\rho>0$ small enough, let
\[
G_\rho:=G\setminus \bigcup_{j=1}^r \overline{D(a_j,\rho)}.
\]
Since $u$ is harmonic on $G_\rho$, Green's identity gives
\[
0=\int_{G_\rho}\Delta u\,dA
  =\int_{\partial G_\rho}\partial_\nu u\,ds
  =\int_{\partial G}\partial_\nu u\,ds
   +\sum_{j=1}^r \int_{\partial D(a_j,\rho)} \partial_\nu u\,ds,
\]
where $\nu$ denotes the outward unit normal of $G_\rho$.

On $\partial G$, the function $u$ decreases in the outward direction, so
\[
\partial_\nu u=-|\nabla u|.
\]
Near $a_j$ one has
\[
u(z)=-m_j\log|z-a_j|+h_j(z),
\]
where $h_j$ is harmonic. Therefore, on $\partial D(a_j,\rho)$,
\[
\partial_\nu u=\frac{m_j}{\rho}+O(1),
\]
because the outward normal of $G_\rho$ on the inner boundary points toward the center $a_j$. Hence
\[
\int_{\partial D(a_j,\rho)} \partial_\nu u\,ds \to 2\pi m_j
\qquad (\rho\downarrow 0).
\]
Letting $\rho\downarrow 0$ yields
\[
-\int_{\partial G}|\nabla u|\,ds + 2\pi\sum_{j=1}^r m_j = 0,
\]
which is exactly the stated identity. \end{proof}
\medskip
\begin{proposition}\label{prop:basic-integral-bound}
\emph{With the notation above,}
\[
I(U)\le 2\pi m.
\]
\end{proposition}

\begin{proof}
Let \(a_1,\dots,a_r\) be the distinct zeros of \(f\) in \(U\), with multiplicities
\(m_1,\dots,m_r\), so that
\[
m_1+\cdots+m_r=m.
\]

We first justify the coarea identities by truncation away from the singularities of
\(u=-\log|f|\). 


Fix regular values $0<\varepsilon<t$, and set
\[
U_\varepsilon:=\{z\in U:u(z)>\varepsilon\},
\qquad
\Omega_{\varepsilon,t}:=U_\varepsilon\setminus K_t=\{\varepsilon<u<t\}.
\]
Since every zero of $f$ in $U$ lies in $K_t$, the function $u$ is harmonic on a neighborhood of $\Omega_{\varepsilon,t}$. Consider
\[
v_{\varepsilon,t}(z):=\min\!\left\{\frac{u(z)-\varepsilon}{t-\varepsilon},\,1\right\}
\qquad (z\in U_\varepsilon).
\]
Then $v_{\varepsilon,t}=0$ on $\partial U_\varepsilon$, $v_{\varepsilon,t}=1$ on $K_t$, and
\[
\nabla v_{\varepsilon,t}=\frac{\nabla u}{t-\varepsilon}
\qquad\text{on }\Omega_{\varepsilon,t}.
\]
Thus $v_{\varepsilon,t}\in H_0^1(U_\varepsilon)$, satisfies $v_{\varepsilon,t}=1$ a.e. on $K_t$, and is therefore admissible in the standard $H^1$-formulation of condenser capacity. By the usual density equivalence between the $C_c^\infty$ and $H_0^1$ admissible classes, this is enough for the capacity computation below.

We now justify directly that $v_{\varepsilon,t}$ minimizes the Dirichlet energy among admissible $H_0^1$-competitors. Let $w\in H_0^1(U_\varepsilon)$ satisfy $w\ge 1$ a.e. on $K_t$. Composing with the $1$-Lipschitz clipping map
\[
\phi(s):=\min\{1,\max\{s,0\}\}
\]
does not increase the Dirichlet energy, so after replacing $w$ by $\phi\circ w$ if necessary, we may assume $0\le w\le 1$ a.e. on $U_\varepsilon$. Since also $w\ge 1$ a.e. on $K_t$, it follows that $w=1$ a.e. on $K_t$.

Because $\varepsilon$ and $t$ are regular values of $u$, the boundaries $\partial U_\varepsilon=\{u=\varepsilon\}$ and $\partial K_t=\{u=t\}$ are smooth, and hence $\Omega_{\varepsilon,t}=U_\varepsilon\setminus K_t$ is a Lipschitz domain with
\[
\partial\Omega_{\varepsilon,t}=\partial U_\varepsilon\cup \partial K_t.
\]
The trace of $w$ on $\partial U_\varepsilon$ is $0$, since $w\in H_0^1(U_\varepsilon)$. Moreover, because $w=1$ a.e. on the interior side $K_t$ and $\partial K_t$ is smooth, the trace of $w|_{\Omega_{\varepsilon,t}}$ on $\partial K_t$ is $1$. The same boundary traces hold for $v_{\varepsilon,t}$: it has trace $0$ on $\partial U_\varepsilon$ and trace $1$ on $\partial K_t$. Therefore
\[
\eta:=w-v_{\varepsilon,t}\in H_0^1(\Omega_{\varepsilon,t}).
\]

Since $u$ is harmonic on a neighborhood of $\Omega_{\varepsilon,t}$, so is $v_{\varepsilon,t}$. Hence, using $\eta$ as a test function in the weak formulation,
\[
\int_{\Omega_{\varepsilon,t}} \nabla v_{\varepsilon,t}\cdot \nabla \eta\, dA = 0.
\]
Expanding $\nabla w=\nabla v_{\varepsilon,t}+\nabla\eta$, we get
\[
\int_{\Omega_{\varepsilon,t}} |\nabla w|^2\,dA
=
\int_{\Omega_{\varepsilon,t}} |\nabla v_{\varepsilon,t}|^2\,dA
+
\int_{\Omega_{\varepsilon,t}} |\nabla \eta|^2\,dA
\ge
\int_{\Omega_{\varepsilon,t}} |\nabla v_{\varepsilon,t}|^2\,dA.
\]
Since $v_{\varepsilon,t}\equiv 1$ on $K_t$, one has $\nabla v_{\varepsilon,t}=0$ a.e. on $K_t$, and therefore
\[
\int_{U_\varepsilon} |\nabla v_{\varepsilon,t}|^2\,dA
=
\int_{\Omega_{\varepsilon,t}} |\nabla v_{\varepsilon,t}|^2\,dA.
\]
Also,
\[
\int_{U_\varepsilon} |\nabla w|^2\,dA
\ge
\int_{\Omega_{\varepsilon,t}} |\nabla w|^2\,dA.
\]
Therefore $v_{\varepsilon,t}$ minimizes the Dirichlet energy among admissible $H_0^1$-competitors, so it is the capacitary potential of $(K_t,U_\varepsilon)$. With the normalization
\[
\operatorname{cap}(E,B)
=
\frac1{4\pi}\inf \int |\nabla v|^2\,dA,
\]
where the infimum is taken over all $v\in H_0^1(B)$ with trace $v\ge 1$ on $E$, we obtain
\[
\operatorname{cap}(K_t,U_\varepsilon)
=
\frac1{4\pi}\int_{\Omega_{\varepsilon,t}} |\nabla v_{\varepsilon,t}|^2\,dA
=
\frac1{4\pi(t-\varepsilon)^2}\int_{\varepsilon<u<t} |\nabla u|^2\,dA.
\]
Applying the coarea formula on the smooth region $\Omega_{\varepsilon,t}$, we obtain
\[
\int_{\varepsilon<u<t} |\nabla u|^2\,dA
=
\int_\varepsilon^t
\left(\int_{u=s} |\nabla u|\,d\sigma\right) ds.
\]
For each regular $s\in(\varepsilon,t)$, the level set $\{u=s\}$ is the disjoint union
of the boundaries of the connected components of $\{u>s\}$. Applying the auxiliary
flux lemma above to each such component and summing over all components shows that
\[
\int_{u=s} |\nabla u|\,d\sigma = 2\pi m,
\]
because the total multiplicity of the zeros of $f$ in $U$ is $m$. Hence
\[
\int_{\varepsilon<u<t} |\nabla u|^2\,dA
=
\int_\varepsilon^t 2\pi m\,ds
=
2\pi m(t-\varepsilon).
\]
Therefore
\[
\operatorname{cap}(K_t,U_\varepsilon)=\frac{m}{2(t-\varepsilon)}.
\]

As $\varepsilon\downarrow 0$, the domains $U_\varepsilon$ increase to $U$. We claim that
\[
\operatorname{cap}_{\mathrm{cond}}(K_t,U)
=
\lim_{\varepsilon\downarrow 0}\operatorname{cap}_{\mathrm{cond}}(K_t,U_\varepsilon).
\]
Indeed, for each $\varepsilon>0$, every admissible function for $(K_t,U_\varepsilon)$
is also admissible for $(K_t,U)$, hence
\[
\operatorname{cap}_{\mathrm{cond}}(K_t,U)
\le
\operatorname{cap}_{\mathrm{cond}}(K_t,U_\varepsilon).
\]
Conversely, let $v\in C_c^\infty(U)$ satisfy $v\ge 1$ on $K_t$. Since $\operatorname{supp}v$
is compact in $U$, there exists $\varepsilon_0>0$ such that $\operatorname{supp}v\subset U_\varepsilon$
for all $0<\varepsilon<\varepsilon_0$. Thus, for all such $\varepsilon$,
\[
\operatorname{cap}_{\mathrm{cond}}(K_t,U_\varepsilon)
\le
\frac1{4\pi}\int_{\mathbb C} |\nabla v|^2\,dA.
\]
Taking $\limsup$ as $\varepsilon\downarrow 0$ and then the infimum over all such $v$
gives the reverse inequality. Therefore
\begin{equation}
\tag{5}
\operatorname{cap}_{\mathrm{cond}}(K_t,U)=\frac{m}{2t}.
\end{equation}


For bounded open sets $G\subset\mathbb C$ we write
\[
\operatorname{cap}_{\log}(G):=\operatorname{cap}_{\log}(\overline G).
\]
We also record explicitly the normalization used here: for the round condenser
\[
\bigl(D(0,R),\overline{D(0,r)}\bigr)
\qquad (0<r<R),
\]
\[
\operatorname{cap}_{\mathrm{cond}}\bigl(\overline{D(0,r)},D(0,R)\bigr)
=
\frac{1}{2\log(R/r)}.
\]
Accordingly, Carleman's inequality in the present normalization reads
\[
\frac{1}{\operatorname{cap}_{\mathrm{cond}}(E,B)}
\le
\log\frac{\operatorname{Area}(B)}{\operatorname{Area}(E)}
\qquad (E\Subset B),
\]
and P\'olya's inequality reads
\[
\operatorname{Area}(K)\le \pi\,\operatorname{cap}_{\log}(K)^2
\qquad (K\subset\mathbb C \text{ compact}).
\]

Applying Carleman's inequality to the condenser \((U,K_t)\) and using \((5)\), we obtain
\[
\log \frac{\operatorname{Area}(U)}{A(t)} \ge \frac{2t}{m},
\]
that is,
\[
A(t)\le \operatorname{Area}(U)e^{-2t/m}.
\tag{6}
\]

It remains to bound \(\operatorname{Area}(U)\). Since \(f\) is monic, the filled lemniscate
\[
\{z\in\mathbb C: |f(z)|\le 1\}
\]
has logarithmic capacity \(1\). Hence, by monotonicity of logarithmic capacity,
\[
\operatorname{cap}_{\log}(\overline U)\le 1.
\]
By Lemma~\ref{lem:level-set-measure-zero}, \(\partial U\subset \{|f|=1\}\) has planar measure zero, so
\[
\operatorname{Area}(U)=\operatorname{Area}(\overline U).
\]
Applying Pólya's inequality to \(\overline U\) yields
\[
\operatorname{Area}(U)=\operatorname{Area}(\overline U)
\le
\pi\,\operatorname{cap}_{\log}(\overline U)^2
\le \pi.
\tag{7}
\]

Combining \((6)\) and \((7)\), we get
\[
A(t)\le \pi e^{-2t/m}.
\]
For every regular $R>0$, the coarea formula on $U\cap\{u<R\}$ gives
\[
\int_{U\cap\{u<R\}} |\nabla u|\,dA
=
\int_0^R P(t)\,dt.
\]
Since $U\cap\{u<R\}\uparrow U$ as $R\to\infty$, monotone convergence yields
\begin{equation}
I(U)=\int_U |\nabla u|\,dA=\int_0^\infty P(t)\,dt.
\tag{1}
\end{equation}
We now record explicitly the regularity of the distribution function
\[
A(t)=\operatorname{Area}(K_t).
\]
Since $u$ has only finitely many critical values and is smooth away from the zeros of
$f$, the function $A$ is absolutely continuous on every compact interval of $(0,\infty)$.
Moreover, for almost every regular value $t>0$,
\[
-A'(t)=\int_{u=t}\frac{1}{|\nabla u|}\,ds.
\]
Indeed, if $[a,b]$ contains no critical values of $u$, then the coarea formula applied
on $\{a<u<b\}$ gives
\[
A(a)-A(b)
=\int_{\{a<u<b\}}1\,dA
=\int_a^b\left(\int_{u=t}\frac{1}{|\nabla u|}\,ds\right)dt.
\]
Since there are only finitely many critical values, this yields the asserted absolute
continuity and derivative formula for almost every $t>0$.

For almost every regular $t>0$, Cauchy--Schwarz and the flux identity give
\[
P(t)^2
=\left(\int_{u=t}1\,ds\right)^2
\le
\left(\int_{u=t}|\nabla u|\,ds\right)
\left(\int_{u=t}\frac{1}{|\nabla u|}\,ds\right)
=
2\pi m\,(-A'(t)).
\]
Hence
\[
P(t)\le \sqrt{2\pi m\,(-A'(t))}. \tag{8}
\]
Using the elementary inequality
\[
2\sqrt{xy}\le \frac{x}{\lambda}+\lambda y
\]
with \(x=2\pi m\), \(y=-A'(t)\), and \(\lambda= m e^{t/m}\), we obtain
\[
P(t)\le \pi e^{-t/m}+\frac{m}{2}e^{t/m}(-A'(t)).
\]
Integrating and using \((1)\),
\[
I(U)\le \pi\int_0^\infty e^{-t/m}\,dt
+\frac{m}{2}\int_0^\infty e^{t/m}(-A'(t))\,dt.
\]
The first integral equals \(\pi m\). For the second, the bound \(A(t)\le \pi e^{-2t/m}\) implies
\(e^{t/m}A(t)\le \pi e^{-t/m}\to 0\) as \(t\to\infty\), so integration by parts gives
\[
\int_0^\infty e^{t/m}(-A'(t))\,dt
=
A(0)+\frac{1}{m}\int_0^\infty e^{t/m}A(t)\,dt.
\]
Using \((7)\) and \((6)\),
\[
A(0)\le \pi,
\qquad
\frac{1}{m}\int_0^\infty e^{t/m}A(t)\,dt
\le
\frac{\pi}{m}\int_0^\infty e^{-t/m}\,dt
=
\pi.
\]
Therefore
\[
\int_0^\infty e^{t/m}(-A'(t))\,dt \le 2\pi,
\]
and hence
\[
I(U)\le \pi m+\frac{m}{2}(2\pi)=2\pi m.
\]
\end{proof}

\medskip
\begin{corollary}\label{cor:component-integral-bound}
\emph{Let $h$ be a monic polynomial, let $c\ge 0$, and let $v:=-\log|h|$. Let $W$ be a connected component of $\{v>c\}$. If $W$
contains exactly $M$ zeros of $h$, counted with multiplicity, then}
\[
\int_W |\nabla v|\,dA\le 2\pi M.
\]
\emph{Equivalently,}
\[
\frac{1}{2\pi}\int_c^{\infty} P_W(t)\,dt\le M,
\]
\emph{where $P_W(t)=\mathcal H^1(\{z\in W:v(z)=t\})$ for regular $t>c$.}
\end{corollary}

\begin{proof}
Set
\[
\widetilde h:=e^c h,
\qquad
\widetilde v:=-\log|\widetilde h|=v-c.
\]
Then $W$ is a connected component of
\[
\{\widetilde v>0\}=\{|\widetilde h|<1\},
\]
and $W$ contains exactly $M$ zeros of $\widetilde h$.

The proof of Proposition~\ref{prop:basic-integral-bound} applies verbatim to $\widetilde h$ on $W$. The only place where monicity was used there was the bound $\operatorname{Area}(W)\le \pi$, obtained via P\'olya's inequality. Here the leading coefficient of $\widetilde h$ is $e^c$, so the filled lemniscate $\{|\widetilde h|\le 1\}$ has logarithmic capacity $e^{-c/\deg h}\le 1$. Since $W\subset \{|\widetilde h|<1\}$, monotonicity gives
\[
\operatorname{cap}_{\log}(W)=\operatorname{cap}_{\log}(\overline W)\le 1.
\]
By Lemma~\ref{lem:level-set-measure-zero}, the boundary $\partial W\subset \{|\widetilde h|=1\}$ has planar measure zero, so
\[
\operatorname{Area}(W)=\operatorname{Area}(\overline W).
\]
Applying P\'olya's inequality to $\overline W$ yields
\[
\operatorname{Area}(W)=\operatorname{Area}(\overline W)
\le
\pi\,\operatorname{cap}_{\log}(W)^2
\le
\pi.
\]
Therefore the same argument gives
\[
\int_W |\nabla \widetilde v|\,dA\le 2\pi M.
\]
Since $\nabla \widetilde v=\nabla v$, this is the desired estimate.
\end{proof}

\section{The first critical value and the low superlevel sets}

Throughout the rest of the paper, $f$ is assumed squarefree on the component under consideration.
Let $U\subset \Lambda_f$ be a connected component containing exactly $m\ge 2$ zeros of $f$, and write again
\[
 u=-\log|f|\quad\text{on }U.
\]

\medskip
\begin{lemma}\label{lem:first-positive-critical-value}
\emph{The function $u$ has at least one critical point in $U$. Consequently the set of positive
critical values of $u$ is nonempty, and it has a smallest element. Denote that smallest positive critical
value by $\lambda_*$ .}
\end{lemma}

\begin{proof} Choose pairwise disjoint closed disks $E_1,\dots,E_m\Subset U$ centered at the zeros of $f$. Since the
superlevel sets of $u$ shrink to the zeros of $f$ as the level tends to $+\infty$, there exists $T>0$ such that
\[
 K_T:=\{u\ge T\}\subset \bigcup_{j=1}^m \operatorname{int}(E_j).
\]
In particular $u<T$ on $\bigcup_j \partial E_j$, and $K_T$ has at least $m$ connected components, one contained in
each $E_j$.

\indent Assume for contradiction that $u$ has no critical point in $\{0<u\le T\}$. Consider the vector field
\[
 X:=\frac{\nabla u}{|\nabla u|^2}
\]
on $\{0<u\le T\}$; it satisfies $Xu=1$. Let $z\in U\setminus K_T$, and write $s:=u(z)\in(0,T)$. Set
\[
 M_s:=U\cap\{e^{-T}\le |f|\le e^{-s/2}\}.
\]
The set $M_s$ is compact. It is bounded because $M_s\subset \Lambda_f$ and $\Lambda_f$ is bounded. To see that it is closed
in $\mathbb{C}$, let $z_k\in M_s$ with $z_k\to z$. Then continuity of $f$ gives
\[
 e^{-T}\le |f(z)|\le e^{-s/2}<1.
\]
Hence $z\in \Lambda_f$. Since $U$ is a connected component of the open set $\Lambda_f$, it is relatively closed in $\Lambda_f$;
because $z$ is also a limit point of $U$, we obtain $z\in U$. Thus $z\in M_s$. Also $M_s\subset\{s/2\le u\le T\}$, so
$X$ is smooth on a neighborhood of $M_s$. Let $\varphi_\tau(z)$ be the maximal trajectory of $X$ through $z$. Along
this trajectory,
\[
 \frac{d}{d\tau}u(\varphi_\tau(z))=Xu(\varphi_\tau(z))=1,
\]
so
\[
 u(\varphi_\tau(z))=s+\tau
\]
whenever the trajectory is defined. In particular, for $0\le \tau\le T-s$ one has
\[
 s/2\le s+\tau\le T,
\]
so the trajectory stays in the compact set $M_s$. Standard ODE continuation therefore shows that
$\varphi_\tau(z)$ is defined on the whole interval $[0,T-s]$, and $\varphi_{T-s}(z)\in K_T$.

\indent Define
\[
R_T(z)=\begin{cases}
z, & z\in K_T,\\
\varphi_{T-u(z)}(z), & z\in U\setminus K_T,
\end{cases}
\]
and
\[
H:[0,1]\times U\to U,\qquad H(\theta,z)=\begin{cases}
z, & z\in K_T,\\
\varphi_{\theta(T-u(z))}(z), & z\in U\setminus K_T.
\end{cases}
\]
If $z_0\in U\setminus K_T$, then the trajectory segment
\[
\{\varphi_\tau(z_0):0\le \tau\le T-u(z_0)\}
\]
is contained in the compact set $M_{u(z_0)}$, on which $X$ is smooth. Standard continuous dependence for smooth ODEs therefore shows that the map
\[
(\theta,z)\mapsto \varphi_{\theta(T-u(z))}(z)
\]
is continuous near each point of $[0,1]\times (U\setminus K_T)$. At points with $u(z)=T$ the two formulas agree, so $H$ is continuous on all of $[0,1]\times U$. Also
\[
H(0,z)=z,\qquad H(1,z)=R_T(z),\qquad H(\theta,z)=z \text{ for every } z\in K_T.
\]
Thus $R_T:U\to K_T$ is a well-defined strong deformation retraction. In particular $K_T$ is connected, contradicting the previous paragraph.
Therefore $u$ has a critical point in $\{0<u\le T\}$, and the set of positive critical values is nonempty.

\indent Since $u=-\log|f|$ is smooth on $U\setminus Z(f)$, its finite critical points are exactly
the points of $U\setminus Z(f)$ where $f'(z)=0$. Because $f$ is squarefree on $U$, we have
\[
Z(f)\cap Z(f')=\varnothing.
\]
Hence $u$ has only finitely many finite critical points in $U$, and therefore only finitely many
positive critical values. Thus the set of positive critical values has a smallest element, denoted
$\lambda_*$. \end{proof}

\medskip
\noindent Set
\[
 \alpha:=\lambda_*/4.
\]
Then $u$ has no critical values in $[0,4\alpha)$.

\medskip
\begin{lemma}\label{lem:superlevel-connected}
\emph{For every $t\in [0,\lambda_*)$, the superlevel set $K_t:=\{u\ge t\}$ is connected. In particular,}
\[
 K:=\{u\ge \alpha/2\},\qquad L:=\{u\ge 3\alpha/2\}
\]
\emph{are compact and connected.}
\end{lemma}

\begin{proof} Fix $0<t<\lambda_*$. By definition of $\lambda_*$, the function $u$ has no critical points in $\{0<u\le t\}$. For
$z\in U\setminus K_t$ write $s:=u(z)\in(0,t)$, and set
\[
 M_s:=U\cap\{e^{-t}\le |f|\le e^{-s/2}\}.
\]
As in the proof of Lemma~\ref{lem:first-positive-critical-value}, $M_s$ is compact and contained in $\{s/2\le u\le t\}$, so the vector field
\[
 X=\frac{\nabla u}{|\nabla u|^2}
\]
is smooth on a neighborhood of $M_s$. Along the maximal trajectory $\varphi_\tau(z)$ of $X$ through $z$ one has
$u(\varphi_\tau(z))=s+\tau$, hence the trajectory stays in $M_s$ for $0\le \tau\le t-s$ and therefore exists on that
whole interval. Consequently
\[
 R_t(z)=\begin{cases}
 z, & z\in K_t,\\
 \varphi_{t-u(z)}(z), & z\in U\setminus K_t,
 \end{cases}
\]
is well defined. Define
\[
H:[0,1]\times U\to U,
\qquad
H(\sigma,z):=\begin{cases}
z, & z\in K_t,\\
\varphi_{\sigma(t-u(z))}(z), & z\in U\setminus K_t.
\end{cases}
\]
As in Lemma~\ref{lem:first-positive-critical-value}, the trajectory segments remain in the compact slab $M_s$, so standard
continuous dependence for smooth ODEs implies that $H$ is continuous. Thus $H$ is a strong
deformation retraction of $U$ onto $K_t$. Since $U$ is connected, $K_t$ is connected. The
case $t=0$ is immediate.

\indent For any $t>0$,
\[
 K_t=U\cap \{|f|\le e^{-t}\},
\]
because $|f|=1$ on $\partial U$. The right-hand side is closed and bounded in $\mathbb{C}$, hence compact. Taking
$t=\alpha/2$ and $t=3\alpha/2$ gives the compactness of $K$ and $L$. \end{proof}

\medskip
\noindent Define
\[
 q:=\frac{1}{2\pi}\int_\alpha^{2\alpha} P(t)\,dt.
\]
Since $P(t)>0$ for $\alpha<t<2\alpha$, we have $q>0$. By Corollary~\ref{cor:component-integral-bound},
\begin{equation}
\frac{1}{2\pi}\int_{2\alpha}^{\infty} P(t)\,dt\le m-q.
\tag{9}
\end{equation}

\section{A generic spanning tree above the base level}

\begin{lemma}[Connected regular regions are strips]\label{lem:regular-strip}
\emph{Let $0<a<b$, and let $V\subset U$ be a connected
open set such that:}
\begin{enumerate}
\item[(i)] \emph{$a<u<b$ on $V$;}
\item[(ii)] \emph{$u$ has no critical points on $V$;}
\item[(iii)] \emph{every component of
\[
 \partial V\setminus \bigl(\{u=a\}\cup\{u=b\}\bigr)
\]
is contained in an integral curve of
\[
 X:=\frac{\nabla u}{|\nabla u|^2}.
\]}
\end{enumerate}
\emph{For $t\in(a,b)$ set
\[
 \Sigma_t:=V\cap\{u=t\}.
\]
Then for every $s\in(a,b)$ the flow of $X$ defines a diffeomorphism
\[
 \Phi_s:\Sigma_s\times(a,b)\to V,\qquad \Phi_s(x,t)=\varphi_{t-s}(x),
\]
where $\varphi_\tau$ denotes the time-$\tau$ flow of $X$. In particular:}
\begin{enumerate}
\item[(a)] \emph{every level section $\Sigma_t$ is connected;}
\item[(b)] \emph{$\partial V\cap\{u=a\}$ is connected;}
\item[(c)] \emph{$\partial V\cap\{u=b\}$ is connected.}
\end{enumerate}
\emph{Thus $V$ is a genuine strip with one lower cross-section and one upper cross-section.}
\end{lemma}

\begin{proof} Since $Xu=1$, the function $u$ increases at unit speed along each trajectory of $X$. Fix $s\in(a,b)$
and $x\in\Sigma_s$. Let $\gamma_x(\tau)=\varphi_\tau(x)$ be the maximal trajectory of $X$ through $x$.

\indent We first claim that $\gamma_x(\tau)$ is defined for all $\tau\in(a-s,b-s)$ and remains in $V$. Fix
\[
 0<\eta<\min\{s-a,b-s\},
\]
and let $K_\eta$ be the closure in $\mathbb{C}$ of the set $V\cap\{a+\eta<u<b-\eta\}$. The set $K_\eta$ is compact. Indeed,
since $0<a<b$ and $0<\eta<\min\{s-a,b-s\}$, we have $a+\eta>0$, so
\[
V\cap\{a+\eta<u<b-\eta\}\subset \{e^{-(b-\eta)}\le |f|\le e^{-(a+\eta)}\},
\]
and the set on the right is closed and bounded in $\mathbb{C}$ because $\Lambda_f$ is bounded. Hence $K_\eta$, being
the closure of a subset of a compact set, is compact. Every point of $K_\eta$ lies in the interior of $U$:
interior points of $V$ are obvious, while points of $K_\eta\cap\partial V$ lie on the side boundary of $V$ and satisfy
$a+\eta\le u\le b-\eta$, so they cannot lie on $\partial U=\{u=0\}$; hence they also lie in the interior of $U$.
By hypothesis each side-boundary component is an integral curve of $X$, so $\nabla u\neq 0$ there as well.
Therefore $X$ is smooth on a neighborhood of $K_\eta$.

\indent Along the trajectory one has
\[
 u(\gamma_x(\tau))=s+\tau
\]
whenever $\gamma_x(\tau)$ is defined. Thus, as long as
\[
 a+\eta\le u(\gamma_x(\tau))\le b-\eta,
\]
the trajectory stays in $K_\eta$. Standard ODE continuation then implies that the trajectory cannot
cease to exist while $u(\gamma_x(\tau))$ remains in $[a+\eta,b-\eta]$.

\indent The only remaining way in which the trajectory could fail to stay in $V$ would be to hit the side
boundary
\[
\partial V\setminus \bigl(\{u=a\}\cup\{u=b\}\bigr)
\]
at some first time $\tau_0\in(a-s,b-s)$. Let $p=\gamma_x(\tau_0)$. By hypothesis, the side-boundary component
through $p$ is itself an integral curve of $X$. Since both that boundary curve and $\gamma_x$ solve the same ODE
and pass through $p$, uniqueness for the initial-value problem forces them to agree in a neighborhood
of $\tau_0$. But the boundary curve lies in $\partial V$, whereas $\gamma_x(\tau)\in V$ for $\tau<\tau_0$ close to $\tau_0$,
a contradiction. Hence $\gamma_x(\tau)$ never hits the side boundary.

\indent Since $u(\gamma_x(\tau))=s+\tau$, the preceding paragraph shows that $\gamma_x(\tau)$ exists and lies in $V$
for every $\tau\in(a-s,b-s)$.

\indent Therefore the formula
\[
\Phi_s(x,t)=\varphi_{t-s}(x)
\]
is well defined on $\Sigma_s\times(a,b)$. It is smooth. If $y\in V$, then the same compact-trapping and first-hit
argument, with $u(y)$ in place of $s$, shows that the trajectory through $y$ exists on the interval
$[s-u(y),0]$ and stays in $V$. Hence $\varphi_{s-u(y)}(y)\in\Sigma_s$, and the inverse map is
\[
 y\mapsto \bigl(\varphi_{s-u(y)}(y),u(y)\bigr),
\]
which is smooth as well. So $\Phi_s$ is a diffeomorphism.

\indent Since $V\cong \Sigma_s\times(a,b)$ and $V$ is connected, the fiber $\Sigma_s$ must be connected. This proves (a). For
$\varepsilon>0$ define
\[
V_\varepsilon^-:=V\cap\{a<u<a+\varepsilon\},\qquad C_\varepsilon^-:=\overline{V_\varepsilon^-},
\]
where the closure is taken in $\mathbb{C}$. By the product representation, $V_\varepsilon^-=\Phi_s(\Sigma_s\times(a,a+\varepsilon))$ is connected,
so $C_\varepsilon^-$ is connected. Since $V$ is bounded, each $C_\varepsilon^-$ is compact. Moreover the family is nested:
\[
 C_{\varepsilon_2}^-\subset C_{\varepsilon_1}^-\qquad (0<\varepsilon_2<\varepsilon_1),
\]
and
\[
 \bigcap_{\varepsilon>0} C_\varepsilon^-=\partial V\cap\{u=a\}.
\]
Hence $\partial V\cap\{u=a\}$ is connected as an intersection of nested compact connected sets. The proof of
(c) is identical, using
\[
V_\varepsilon^+:=V\cap\{b-\varepsilon<u<b\}.
\]
\end{proof}

\medskip
\begin{lemma}[Regular lemniscate components are Jordan domains]\label{lem:jordan-domain}
\emph{Let $c$ be a regular value of
$u=-\log|f|$, and let $W$ be a connected component of
\[
\{u>c\}=\{|f|<e^{-c}\}.
\]
Then $W$ is simply connected, and $\partial W$ is a smooth Jordan curve.}
\end{lemma}

\begin{proof} Since $c$ is a regular value of $u$, the set $\partial W\subset\{u=c\}$ is a smooth compact $1$-manifold, hence
a finite disjoint union of smooth Jordan curves.

\indent Suppose that $W$ had more than one boundary component. Choose an innermost boundary
component $\Gamma\subset\partial W$ in the following sense: if $H$ is the bounded Jordan domain enclosed by $\Gamma$, then
$H$ contains no other boundary component of any connected component of $\{u>c\}$. Such a choice is
possible because the boundary consists of finitely many disjoint Jordan curves.

\indent By the choice of $\Gamma$, the interior $H$ contains no point of $\{u>c\}$. Otherwise some connected
component of $\{u>c\}$ would lie in $H$, and one of its boundary components would be strictly inside
$\Gamma$, contradicting innermostness. Hence
\[
H\subset\{u\le c\}=\{|f|\ge e^{-c}\},\qquad \partial H=\Gamma\subset\{u=c\}.
\]
In particular $f$ has no zeros on $H$, so
\[
 v:=\log|f|
\]
is harmonic on $H$ and continuous on $\overline H$. On $\partial H$ we have $v\equiv -c$. By the maximum principle,
$v\equiv -c$ on $H$. Thus $|f|$ is constant on $H$, impossible for a nonconstant holomorphic function.

\indent This contradiction shows that $\partial W$ has only one connected component. Since it is a smooth
compact $1$-manifold, it is a single smooth Jordan curve. Therefore $W$ is simply connected. \end{proof}

\medskip
\begin{lemma}[Regular-slab deformation and connectivity]\label{lem:regular-slab}
\emph{Let $V$ be a connected component of
$\{u>c\}$, and let $c<a<b$. Assume that $u$ has no critical points on
\[
 M:=V\cap\{a\le u\le b\}.
\]
Write
\[
V_{\ge t}:=V\cap\{u\ge t\},\qquad V_{>t}:=V\cap\{u>t\}.
\]
Then:}
\begin{enumerate}
\item[(a)] \emph{the map
\[
R_{a,b}:V_{\ge a}\to V_{\ge b},\qquad R_{a,b}(z)=\begin{cases}
z, & u(z)\ge b,\\
\varphi_{b-u(z)}(z), & a\le u(z)<b,
\end{cases}
\]
is a well-defined strong deformation retraction;}
\item[(b)] \emph{if $V_{\ge b}$ is connected, then $V_{>a}$ is connected. More precisely, every point of $V_{>a}$ can be joined to
$V_{\ge b}$ by a trajectory segment of $X=\nabla u/|\nabla u|^2$ lying entirely in $V_{>a}$;}
\item[(c)] \emph{if moreover $a$ and $b$ are regular values, then every connected component of $V\cap\{a<u<b\}$ is
 a regular strip in the sense of Lemma~\ref{lem:regular-strip}.}
\end{enumerate}
\end{lemma}

\begin{proof} Since $u>c$ on $V$ and $a>c$, the compact slab $M$ is contained in the interior of $V$. On $M$
the vector field
\[
 X:=\frac{\nabla u}{|\nabla u|^2}
\]
is smooth and satisfies $Xu=1$.

\indent If $z\in V_{\ge a}$ and $u(z)<b$, then along the trajectory of $X$ through $z$ the value of $u$ increases
linearly:
\[
 u(\varphi_\tau(z))=u(z)+\tau.
\]
Hence the trajectory is defined for all $0\le \tau\le b-u(z)$ and remains in $M\subset V$. This proves that
$R_{a,b}$ is well defined. Define
\[
 H:[0,1]\times V_{\ge a}\to V_{\ge a}
\]
by
\[
H(s,z)=\begin{cases}
z, & u(z)\ge b,\\
\varphi_{s(b-u(z))}(z), & a\le u(z)<b.
\end{cases}
\]
If $z_0\in V_{\ge a}$ and $u(z_0)<b$, then the trajectory segment
\[
\{\varphi_\tau(z_0):0\le \tau\le b-u(z_0)\}
\]
lies in the compact slab $M$, on which $X$ is smooth. Standard continuous dependence for smooth ODEs therefore implies that $H$ is continuous near each point of $[0,1]\times (V_{\ge a}\setminus V_{\ge b})$. On $V_{\ge b}$ we have $H(s,z)=z$, and at points with $u(z)=b$ the two formulas agree. Therefore $H$ is continuous on all of $[0,1]\times V_{\ge a}$. It is immediate that
\[
H(0,z)=z,\qquad H(1,z)=R_{a,b}(z),\qquad H(s,z)=z\text{ for }z\in V_{\ge b}.
\]
Thus $R_{a,b}$ is a strong deformation retraction. This proves (a).

\indent For (b), let $z\in V_{>a}$. If $u(z)\ge b$, then $z\in V_{\ge b}$. Otherwise the trajectory segment
\[
\tau\mapsto \varphi_\tau(z),\qquad 0\le \tau\le b-u(z),
\]
lies in $V_{>a}$ and ends at $R_{a,b}(z)\in V_{\ge b}$. Therefore every point of $V_{>a}$ lies in the same connected
component of $V_{>a}$ as some point of $V_{\ge b}$. If $V_{\ge b}$ is connected, it follows that $V_{>a}$ is connected.

\indent For (c), let $S$ be a connected component of $V\cap\{a<u<b\}$. Since $\partial V\subset\{u=c\}$ and $c<a$,
every boundary point of $S$ lying in $V$ belongs to one of the levels $\{u=a\}$ or $\{u=b\}$. Thus
\[
\partial S\setminus \bigl(\{u=a\}\cup\{u=b\}\bigr)=\varnothing,
\]
so the side-boundary hypothesis of Lemma~\ref{lem:regular-strip} is vacuous. Applying Lemma~\ref{lem:regular-strip} to $S$ gives the desired
strip structure. \end{proof}


\paragraph{Component inheritance for nested superlevel sets.}
We shall use the following elementary observation repeatedly. Let \(V\) be a connected component of
\[
\{u>c\},
\]
and let \(d>c\). If \(C\) is a connected component of
\[
V\cap\{u>d\},
\]
then \(C\) is also a connected component of the global superlevel set \(\{u>d\}\). Indeed, if \(D\) is the
connected component of \(\{u>d\}\) containing \(C\), then \(D\subset \{u>c\}\). Since \(D\) is connected
and meets the connected component \(V\) of \(\{u>c\}\), we must have \(D\subset V\). Hence
\(D\subset V\cap\{u>d\}\), and by maximality of \(C\) inside \(V\cap\{u>d\}\) we get \(D=C\).

\medskip
\begin{lemma}[Geometric realization of a finite tree]\label{lem:geometric-tree-realization}
\emph{Let $T$ be a finite tree with vertex set $V(T)$ and
edge set $E(T)$. Suppose that to each vertex $v\in V(T)$ there is assigned a point $p_v\in U$, and to each
edge $e=\{v,w\}\in E(T)$ there is assigned an embedded arc
\[
 A_e\subset U
\]
joining $p_v$ to $p_w$, such that:}
\begin{enumerate}
\item[(i)] \emph{if $v\neq w$, then $p_v\neq p_w$;}
\item[(ii)] \emph{the interior of $A_e$ contains no point $p_{v'}$;}
\item[(iii)] \emph{if $e\neq e'$, then the interiors of $A_e$ and $A_{e'}$ are disjoint;}
\item[(iv)] \emph{if $e=\{v,w\}$ and $e'=\{v',w'\}$ are distinct, then
\[
 A_e\cap A_{e'}\subset \{p_v,p_w\}\cap\{p_{v'},p_{w'}\}.
\]
Equivalently, two distinct edge-arcs meet only at common endpoints, and they meet there
exactly when the corresponding edges of $T$ share that vertex.}
\end{enumerate}
\emph{Then
\[
 H:=\bigcup_{e\in E(T)} A_e
\]
is homeomorphic to $T$.}
\end{lemma}

\begin{proof} For each edge $e=\{v,w\}$ choose a homeomorphism $\varphi_e:[0,1]\to A_e$ with $\varphi_e(0)=p_v$ and
$\varphi_e(1)=p_w$. Because of (ii), no interior point of an edge-arc is identified with a vertex-point. Because
of (iii) and (iv), if $e\neq e'$ then points of $[0,1]$ and $[0,1]$ can map to the same point of $U$ only when
both are endpoints corresponding to a common endpoint-vertex of the two edges. Hence the maps
$\varphi_e$ are compatible with the quotient relation defining the geometric realization $|T|$ of the abstract
tree. They therefore glue to a continuous map
\[
 \Phi:|T|\to H.
\]
The map $\Phi$ is surjective by construction. It is injective because distinct open edges of $|T|$ map
into disjoint interiors by (iii), and an open edge can meet another edge-image only at shared
endpoint-vertices by (iv). Thus $\Phi$ is a continuous bijection from the compact space $|T|$ onto the
Hausdorff space $H\subset\mathbb{C}$, hence a homeomorphism. \end{proof}

\medskip
\noindent\textbf{Auxiliary topological fact.}
Let $\Omega_0,\Omega_1\subset \mathbb C$ be Jordan domains such that
\[
\Omega_1\subset \Omega_0
\qquad\text{and}\qquad
\partial\Omega_1\cap \partial\Omega_0=\varnothing.
\]
Then $\overline{\Omega_1}\subset \Omega_0$, and $\Omega_0\setminus \overline{\Omega_1}$
is homeomorphic to an annulus.

\smallskip
\noindent
Indeed, since $\overline{\Omega_1}$ is connected, meets $\Omega_0$, and is disjoint from
$\partial\Omega_0$, it cannot meet the exterior component of $\mathbb C\setminus \partial\Omega_0$.
Hence $\overline{\Omega_1}\subset \Omega_0$.
Now choose, by the Jordan--Schoenflies theorem, a homeomorphism
\[
h:\overline{\Omega_0}\to \overline{\mathbb D}.
\]
Then $h(\partial\Omega_1)$ is a Jordan curve contained in $\mathbb D$, so it bounds a Jordan
domain $\Delta\Subset \mathbb D$. The set $h(\Omega_1)$ is a connected component of
$\mathbb D\setminus h(\partial\Omega_1)$, whose two components are $\Delta$ and
$\mathbb D\setminus \overline{\Delta}$. Since $h(\overline{\Omega_1})$ is compactly contained in
$\mathbb D$, its closure does not meet $\partial\mathbb D$, so $h(\Omega_1)$ cannot be the outer
component $\mathbb D\setminus \overline{\Delta}$. Hence $h(\Omega_1)=\Delta$. Therefore
\[
h\bigl(\Omega_0\setminus \overline{\Omega_1}\bigr)=\mathbb D\setminus \overline{\Delta},
\]
which is an annulus.

\medskip
\begin{proposition}\label{prop:generic-spanning-tree}
\emph{Let $0\le c<2\alpha$, and let $V$ be a connected component of $\{u>c\}$ containing exactly $m$
zeros of $f$. Assume that:}
\begin{enumerate}
\item[(a)] \emph{$u$ has no critical values on $V$ in $[c,2\alpha]$;}
\item[(b)] \emph{all critical points of $u$ in $V$ with value $>2\alpha$ are nondegenerate;}
\item[(c)] \emph{all critical values of $u$ in $V$ larger than $2\alpha$ are pairwise distinct.}
\end{enumerate}
\emph{Then for every $\varepsilon>0$ there exists a finite embedded tree $G_\varepsilon\subset V$ containing all $m$ zeros of $f$ in $V$
and satisfying}
\begin{equation}
\operatorname{len}(G_\varepsilon)\le \frac{1}{2\pi}\int_{2\alpha}^{\infty} P_V(t)\,dt+\varepsilon,
\tag{10}
\end{equation}
\emph{where $P_V(t):=\mathcal H^1(\{z\in V:u(z)=t\})$ for regular $t>2\alpha$.}
\end{proposition}

\begin{proof}

If $m=1$, then the one-point set consisting of the unique zero has length $0$, so the conclusion is immediate. Thus we may assume that $m \ge 2$.

Let the distinct critical values of $u$ in $V$ above $2\alpha$ be
\[
2\alpha < \mu_1 < \mu_2 < \cdots < \mu_s.
\]
We claim that $s \ge 1$.

Indeed, suppose for contradiction that $s=0$. Then there are no critical values of $u$ in $(2\alpha,\infty)$. We first claim that
\[
V_{\ge 2\alpha}:=\{z\in V : u(z)\ge 2\alpha\}
\]
is connected.

Let $z\in V$ with $u(z)<2\alpha$, and write
\[
s_z:=u(z)\in(c,2\alpha).
\]
Set
\[
M_z:=V\cap\{s_z\le u\le 2\alpha\}.
\]
Because $V$ is a connected component of the open set $\{u>c\}$, it is relatively closed
in $\{u>c\}$; since $s_z>c$, the set $M_z$ is closed in $\mathbb C$. It is also bounded
because $\Lambda_f$ is bounded, hence compact.

Moreover, $u$ has no critical points on $M_z$, so the vector field
\[
X:=\frac{\nabla u}{|\nabla u|^2}
\]
is smooth on a neighborhood of $M_z$. Let $\phi_\tau(z)$ be the maximal trajectory of
$X$ through $z$. Along this trajectory one has
\[
\frac{d}{d\tau}u(\phi_\tau(z))
= \nabla u(\phi_\tau(z)) \cdot X(\phi_\tau(z))
= 1,
\]
hence, as long as the trajectory is defined,
\[
u(\phi_\tau(z)) = s_z + \tau.
\]
Therefore, for every $0 \le \tau \le 2\alpha - s_z$, the point $\phi_\tau(z)$ belongs to
$M_z$. Since $X$ is smooth on a neighborhood of the compact set $M_z$, the standard
continuation theorem for ODEs implies that the trajectory extends throughout the whole
interval $[0,2\alpha-s_z]$. In particular,
\[
\phi_{2\alpha-s_z}(z) \in V_{\ge 2\alpha}.
\]

Define
\[
H : [0,1]\times V \to V,
\qquad
H(\theta,z) :=
\begin{cases}
\phi_{\theta(2\alpha-u(z))}(z), & u(z)<2\alpha,\\[1mm]
z, & u(z)\ge 2\alpha.
\end{cases}
\]
We claim that $H$ is continuous. Fix $s$ with $c<s<2\alpha$. On the compact slab
\[
V \cap \{s \le u \le 2\alpha\},
\]
the vector field $X$ is smooth on a neighborhood, so its flow depends continuously on
initial data and time. Hence $H$ is continuous on
\[
[0,1]\times \bigl(V\cap\{u\ge s\}\bigr).
\]
Since $s\in(c,2\alpha)$ is arbitrary, these local continuity statements patch together to
show that $H$ is continuous on $[0,1]\times V$.

Thus $H$ is a strong deformation retraction of $V$ onto $V_{\ge 2\alpha}$. Since $V$ is
connected, $V_{\ge 2\alpha}$ is connected.

Now let $T>2\alpha$ be a sufficiently large regular value such that $\{z\in V:u(z)>T\}$ has exactly $m$ connected components, one near each zero of $f$ in $V$. Since there are no critical values of $u$ in $[2\alpha,T]$, Lemma~\ref{lem:regular-slab}(a) gives a strong deformation retraction of $V_{\ge 2\alpha}$ onto
\[
V_{\ge T}:=\{z\in V:u(z)\ge T\}.
\]
Hence $V_{\ge T}$ is connected.

Let $W_1,\dots,W_m$ be the connected components of $\{z\in V:u(z)>T\}$. Since $T$ is a regular value of $u$, every point $p\in\{z\in V:u(z)=T\}$ has a neighborhood $U_p\subset V$ in which, after a smooth change of coordinates, $u-T=x$. Hence
\[
U_p\cap\{u>T\}
\]
is connected. It follows that at most one component $W_j$ can accumulate at $p$. Therefore the closures $\overline{W_1}^{\,V},\dots,\overline{W_m}^{\,V}$ are pairwise disjoint. Moreover,
\[
V_{\ge T}=\bigcup_{j=1}^m \overline{W_j}^{\,V}.
\]
Since the closure of a connected set is connected, each $\overline{W_j}^{\,V}$ is connected. Thus these $m$ sets are precisely the connected components of $V_{\ge T}$. Because $m\ge2$, this contradicts the connectedness of $V_{\ge T}$.

This contradiction shows that $s\ge1$.

\indent Choose pairwise disjoint closed intervals
\[
 I_r=[\mu_r-\delta_r,\mu_r+\delta_r] \qquad (1\le r\le s)
\]
so small that each $I_r$ contains no critical value of $u$ other than $\mu_r$, and so small that the local Morse neighborhoods constructed below can be chosen with associated length bounds $\eta_r>0$ satisfying
\[
\sum_{r=1}^s \eta_r<\frac{\varepsilon}{2}.
\]
Choose a regular value
\[
 \tau_0\in(2\alpha,\mu_1-\delta_1).
\]

For each zero $a_j$, write
\[
f(z)=(z-a_j)^{m_j}h_j(z)
\]
on a neighborhood $U_j$ of $a_j$, where $h_j(a_j)\neq 0$. Choose pairwise disjoint closed disks
\[
B(a_j,r_j)\Subset U_j
\]
and a holomorphic branch $q_j$ of $h_j^{1/m_j}$ on $B(a_j,r_j)$. Set
\[
\psi_j(z):=(z-a_j)q_j(z).
\]
After shrinking $r_j$ if necessary, we may assume that $\psi_j$ is biholomorphic on a neighborhood of $B(a_j,r_j)$ and that
\[
f(z)=\psi_j(z)^{m_j}
\qquad (z\in B(a_j,r_j)).
\]
Since $\psi_j^{-1}$ is $C^1$ near $0$, there exist constants $C_j>0$ and $\rho_j>0$ such that, for every $w$ with $|w|\le \rho_j$, the preimage of the radial segment $[0,w]$ under $\psi_j^{-1}$ has length at most $C_j|w|$.

Choose a regular value
\[
 T>\mu_s+\delta_s
\]
so large that the connected components
\[
 Z_1,\dots,Z_m
\]
of $\{z\in V:u(z)>T\}$ have pairwise disjoint closures and, in addition, satisfy
\begin{enumerate}
\item[(i)] $\overline{Z_j}\subset B(a_j,r_j)$ for every $j$;
\item[(ii)] $e^{-T/m_j}\le \rho_j$ for every $j$;
\item[(iii)] $\sum_{j=1}^m C_j e^{-T/m_j}<\varepsilon/2$.
\end{enumerate}

For each $j$ set
\[
K_j:=\overline{Z_j}.
\]
Then $K_1,\dots,K_m$ are pairwise disjoint compact sets, each lying in a small neighborhood of one zero
$a_j$ of $f$. Since $f$ is squarefree, each $Z_j$ (and hence each $K_j$) contains exactly one zero $a_j$.

\indent We first note that $\{z\in V:u>\tau_0\}$ is connected. Choose a regular value $t_0\in(c,2\alpha)$. As above,
the absence of critical points on $V\cap\{c<u\le t_0\}$ gives a strong deformation retraction of $V$ onto
\[
V_{\ge t_0}:=\{z\in V:u\ge t_0\},
\]
so $V_{\ge t_0}$ is connected. Since there are no critical points in $\{z\in V:t_0\le u\le \tau_0\}$, Lemma~\ref{lem:regular-slab}(a) shows
that
\[
V_{\ge \tau_0}:=\{z\in V:u\ge \tau_0\}
\]
is connected. Next choose a regular value $\tau_1\in(\tau_0,\mu_1-\delta_1)$. Applying Lemma~\ref{lem:regular-slab}(a) on $[\tau_0,\tau_1]$,
we find that $V_{\ge \tau_1}$ is connected. Lemma~\ref{lem:regular-slab}(b) applied with $a=\tau_0$ and $b=\tau_1$ now gives that
$\{z\in V:u>\tau_0\}$ is connected.

\indent For each $r=1,\dots,s$, let $p_r$ be the unique critical point of $u$ with critical value $\mu_r$. Since $p_r$ is
nondegenerate and $u$ is harmonic, the Morse normal form gives local coordinates $(x,y)$ centered at
$p_r$ in which
\[
 u=\mu_r+x^2-y^2;
\]
see, for example, \cite{Milnor1963}. Choose pairwise disjoint closed neighborhoods $N_r$ of $p_r$ so small that:
\begin{enumerate}
\item[(i)] $N_r\subset u^{-1}(I_r)$;
\item[(ii)] $\partial N_r$ is the union of one connected incoming cross-section
\[
 C_r^-\subset\{u=\mu_r-\delta_r\},
\]
two connected outgoing cross-sections
\[
 C_{r,1}^+,\ C_{r,2}^+\subset\{u=\mu_r+\delta_r\},
\]
and finitely many flow segments of $X$ joining the endpoints of these cross-sections;
\item[(iii)] inside $N_r$, the local superlevel picture is the standard saddle picture: crossing the level $\mu_r$
      upward changes one component into two components;
\item[(iv)] outside $N_r$, there are no critical points in the slab
\[
 \mu_r-\delta_r\le u\le \mu_r+\delta_r;
\]
\item[(v)] there exists a number $\eta_r>0$ such that, for every choice of points
\[
x_{r,-}\in C_r^-,
\qquad
x_{r,1}\in C_{r,1}^+,
\qquad
x_{r,2}\in C_{r,2}^+,
\]
there is an embedded $Y$-shaped tree
\[
T_r(x_{r,-},x_{r,1},x_{r,2})\subset N_r
\]
with branch point $p_r$ and terminal vertices $x_{r,-},x_{r,1},x_{r,2}$, and
\[
\operatorname{len}\bigl(T_r(x_{r,-},x_{r,1},x_{r,2})\bigr)\le \eta_r;
\]
\item[(vi)] there exist connected open collar neighborhoods
\[
\mathcal C_r^-,
\quad
\mathcal C_{r,1}^+,
\quad
\mathcal C_{r,2}^+
\subset (U_r\setminus N_r)\cap\{\tau_0<u<T\}
\]
such that
\[
C_r^- \subset \overline{\mathcal C_r^-},
\qquad
C_{r,i}^+ \subset \overline{\mathcal C_{r,i}^+} \quad (i=1,2),
\]
and every point of $(U_r\setminus N_r)\cap\{\tau_0<u<T\}$ sufficiently close to $C_r^-$
(respectively to $C_{r,i}^+$) on the exterior side of $N_r$ lies in $\mathcal C_r^-$
(respectively in $\mathcal C_{r,i}^+$).
\end{enumerate}
Such neighborhoods exist. Fix $r$ and work in Morse coordinates $(x,y)$ centered at $p_r$ with
\[
u=\mu_r+x^2-y^2.
\]
Take a sufficiently small standard saddle neighborhood whose boundary consists of one arc on the level $\{u=\mu_r-\delta_r\}$, two arcs on the level $\{u=\mu_r+\delta_r\}$, and four short flow segments, arranged in the usual local saddle picture. For each of the three level arcs, let $q_{r,-},q_{r,1},q_{r,2}$ denote the endpoints of the local separatrix segments meeting those arcs. If
\[
x_{r,-}\in C_r^-,
\qquad
x_{r,1}\in C_{r,1}^+,
\qquad
x_{r,2}\in C_{r,2}^+,
\]
then each $x_{r,*}$ can be joined inside its cross-section to the corresponding $q_{r,*}$ by a boundary subarc, and the union of these three subarcs with the three separatrix segments from $q_{r,-},q_{r,1},q_{r,2}$ to $p_r$ is an embedded $Y$-shaped tree in $N_r$ with branch point $p_r$ and terminal vertices $x_{r,-},x_{r,1},x_{r,2}$. As the neighborhood is shrunk, both the cross-sections and the separatrix segments shrink to $p_r$, so the total length of this tree tends to $0$ uniformly in the choices of $x_{r,-},x_{r,1},x_{r,2}$; this gives property \emph{(v)} for some $\eta_r>0$. Likewise, the exterior saddle sectors provide the one-sided connected collar neighborhoods required in \emph{(vi)}. Since there are only finitely many critical points, we may choose the neighborhoods $N_r$ small enough that, in addition,
\[
\sum_{r=1}^s \eta_r<\frac{\varepsilon}{2}.
\]
Set
\[
\Omega:=\{z\in V:\tau_0<u<T\}\setminus \bigcup_{r=1}^s N_r.
\]

We shall prove below that, for the neighborhoods $N_r$ chosen above, each of the cross-sections $C_r^{-}, C_{r,1}^{+}, C_{r,2}^{+}$
lies on the boundary of a unique connected component of $\Omega$.

Let $S$ be a connected component of $\Omega$. Since $u$ has no critical points on $S$ and every side-boundary
component of $S$ is a flow line of $X$, Lemma~\ref{lem:regular-strip} applies: $S$ is a regular strip, with one connected lower
cross-section and one connected upper cross-section, and every level section
\[
\Gamma_t:=S\cap\{u=t\}
\]
is connected.

\indent For such a strip $S$, let $a_S<b_S$ be the levels of its lower and upper cross-sections.
For $t\in(a_S,b_S)$, let $C_t(S)$ denote the connected component of
\[
\{z\in V:u(z)>t\}
\]
that meets $S\cap\{t<u<t+\eta\}$ for some (equivalently, every sufficiently small)
$\eta>0$. Let $k_S$ be the number of zeros of $f$ in $C_t(S)$. This is independent of
$t\in(a_S,b_S)$. Indeed, if $t_1<t_2$ both lie in $(a_S,b_S)$, then the slab
\[
M:=C_{t_1}(S)\cap\{t_1\le u\le t_2\}
\]
contains no critical points of $u$. Hence the flow of
\[
X=\frac{\nabla u}{|\nabla u|^2}
\]
provides the regular continuation of the component determined by $S$ from level $t_1$ up to level $t_2$; in particular, no splitting or merging of this component can occur between these two regular levels. Thus the component of $\{u>t_2\}$ obtained from $C_{t_1}(S)$ by passing through the slab is exactly $C_{t_2}(S)$. Since every zero of $f$ corresponds to the singular level $u=+\infty$, no zero lies in the finite slab $\{t_1\le u\le t_2\}$, and therefore the set of zeros contained in the continued component is the same at levels $t_1$ and $t_2$. Hence the number of zeros in $C_t(S)$ is constant on $(a_S,b_S)$.

\begin{lemma}\label{lem:strip-component-contains-zero}
For every strip $S$ and every $t \in (a_S,b_S)$, the component $C_t(S)$ contains at
least one zero of $f$. In particular,
\[
k_S \ge 1.
\]
\end{lemma}

\begin{proof}
Fix $t \in (a_S,b_S)$ and set
\[
W := C_t(S).
\]
Since $t$ is a regular value and $W$ is a connected component of
\[
\{z \in V : u(z) > t\},
\]
we have
\[
\partial W \subset \{u=t\}.
\]
In particular, no zero of $f$ lies on $\partial W$, because $u=+\infty$ at every zero.

Assume for contradiction that $W$ contains no zero of $f$. Then $f \neq 0$ on
$\overline W$, so
\[
u=-\log|f|
\]
is harmonic on an open neighborhood of $\overline W$. Since $u=t$ on $\partial W$,
the maximum principle yields
\[
u \le t \quad \text{throughout } W,
\]
contradicting
\[
W \subset \{u>t\}.
\]
Therefore $W$ contains a zero of $f$, as claimed.
\end{proof}

Set
\[
\Gamma_t:=S\cap\{u=t\}.
\]
We claim that
\[
\partial C_t(S)=\Gamma_t.
\]

Since $t$ is a regular value and $C_t(S)$ is a connected component of $\{u>t\}\cap V$,
Lemma~\ref{lem:jordan-domain} implies that $\partial C_t(S)$ is a smooth Jordan curve contained in
$\{u=t\}\cap V$.

Because $t\in (a_S,b_S)\subset (\tau_0,T)\setminus \bigcup_{r=1}^s I_r$, the level
$\{u=t\}\cap V$ does not meet any $N_r$. Hence
\[
\{u=t\}\cap V=\{u=t\}\cap \Omega,
\]
and $\Gamma_t$ is a connected component of $\Omega\cap\{u=t\}$.

First, $\Gamma_t\subset \partial C_t(S)$: if $x\in \Gamma_t$, then for every sufficiently
small $\eta>0$ the set $S\cap\{t<u<t+\eta\}$ lies in $C_t(S)$ by definition, and its
points converge to $x$. Since $x\notin C_t(S)$, this shows $x\in \partial C_t(S)$.

Conversely, $\partial C_t(S)$ is a connected subset of $\Omega\cap\{u=t\}$, and it meets
$\Gamma_t$ by the previous paragraph. Because $\Gamma_t$ is a connected component of
$\Omega\cap\{u=t\}$, it follows that
\[
\partial C_t(S)\subset \Gamma_t.
\]
Combining the two inclusions gives
\[
\partial C_t(S)=\Gamma_t.
\]
Applying the auxiliary flux lemma to the component $C_t(S)$ gives
\[
\int_{\Gamma_t} |\nabla u|\,ds = 2\pi k_S.
\]
For an arc $I\subset \Gamma_{t_1}$, let $I_{t_2}$ be its image under the flow from level $t_1$ to level $t_2$. The region swept by
these trajectories has boundary consisting of $I$, $I_{t_2}$, and two flow lines. Since $\Delta u=0$ in this region
and $\nabla u$ is tangent to the side flow lines, Green's formula gives
\[
\int_I |\nabla u|\,ds = \int_{I_{t_2}} |\nabla u|\,ds.
\]
Thus the measures
\[
 d\mu_t:=\frac{|\nabla u|\,ds}{2\pi k_S}
\]
on $\Gamma_t$ are probability measures preserved by the flow.

\indent Fix $s_0\in(a_S,b_S)$. For $x\in\Gamma_{s_0}$, let
\[
\widetilde\gamma_x(t):=\Phi_{s_0}(x,t),\qquad t\in(a_S,b_S),
\]
where $\Phi_{s_0}$ is the strip diffeomorphism from Lemma~\ref{lem:regular-strip}. We claim that $\widetilde\gamma_x$ extends continuously to
$[a_S,b_S]$, with endpoint values on the lower and upper cross-sections of $S$. Indeed, let $K_x^-$ and $K_x^+$
be the closures in $\mathbb{C}$ of the sets
\[
 S\cap\{a_S<u<s_0\}, \qquad S\cap\{s_0<u<b_S\},
\]
respectively. These sets are compact because
\[
S\cap\{a_S<u<s_0\}\subset\{a_S\le u\le s_0\},\qquad S\cap\{s_0<u<b_S\}\subset\{s_0\le u\le b_S\},
\]
and the right-hand sides are closed and bounded subsets of $\mathbb{C}$. Moreover every point of $K_x^{\pm}$ is
noncritical: interior points because $S$ contains no critical points, points on the lower or upper
cross-sections because these lie on regular levels of $u$, and points on the side boundary because
those components are integral curves of $X$. All such points lie in the interior of $V\setminus\{p_1,\dots,p_s\}$,
so by continuity $X$ is smooth on a neighborhood of $K_x^-\cup K_x^+$. Since $Xu=1$, standard ODE
continuation shows that the trajectory through $x$ extends to times $a_S-s_0$ and $b_S-s_0$. Thus
\[
\gamma_x:=\widetilde\gamma_x([a_S,b_S])
\]
is a well-defined compact trajectory segment joining the lower cross-section of $S$ to its upper
cross-section. For $x\in\Gamma_{s_0}$, define
\[
\ell(x):=\operatorname{len}(\gamma_x).
\]
Since $\widetilde\gamma_x(t)=\Phi_{s_0}(x,t)$ for $t\in[a_S,b_S]$, we have
\[
\ell(x)
=
\int_{a_S}^{b_S}\left|\partial_t\Phi_{s_0}(x,t)\right|\,dt
=
\int_{a_S}^{b_S}\left|X\bigl(\Phi_{s_0}(x,t)\bigr)\right|\,dt.
\]
Because
\[
X=\frac{\nabla u}{|\nabla u|^2},
\]
it follows that
\[
|X|=\frac{1}{|\nabla u|},
\]
and hence
\[
\ell(x)
=
\int_{a_S}^{b_S}\frac{dt}{|\nabla u(\Phi_{s_0}(x,t))|}.
\]
Integrating with respect to the probability measure $\mu_{s_0}$ and using Fubini,
\[
\int_{\Gamma_{s_0}}\ell(x)\,d\mu_{s_0}(x)
=
\int_{a_S}^{b_S}
\left(
\int_{\Gamma_{s_0}}
\frac{1}{|\nabla u(\Phi_{s_0}(x,t))|}
\,d\mu_{s_0}(x)
\right)\,dt.
\]
Since the flow preserves the measures $\mu_t$, we have
\[
(\Phi_{s_0}(\cdot,t))_{\#}\mu_{s_0}=\mu_t.
\]
Therefore the inner integral equals
\[
\int_{\Gamma_t}\frac{1}{|\nabla u(y)|}\,d\mu_t(y)
=
\frac{1}{2\pi k_S}\int_{\Gamma_t} ds
=
\frac{P_S(t)}{2\pi k_S},
\]
where $P_S(t):=\mathcal H^1(\Gamma_t)$. Thus
\[
\int_{\Gamma_{s_0}} \operatorname{len}(\gamma_x)\,d\mu_{s_0}(x)
=
\frac{1}{2\pi k_S}\int_{a_S}^{b_S} P_S(t)\,dt.
\]
Since $\mu_{s_0}$ is a probability measure, there exists $x\in \Gamma_{s_0}$ such that
\[
\operatorname{len}(\gamma_x)\le \frac{1}{2\pi k_S}\int_{a_S}^{b_S} P_S(t)\,dt.
\]
Fix such an $x$ and denote the corresponding trajectory segment by $\gamma_S$.

\indent Summing over all strips and using Lemma~\ref{lem:strip-component-contains-zero}, we obtain
\[
\sum_S \operatorname{len}(\gamma_S)\le \frac{1}{2\pi}\sum_S\int_{a_S}^{b_S} P_S(t)\,dt.
\]
Because $\Omega$ is an open subset of the plane, it has at most countably many connected components, so the family of strips is at most countable. Since every summand is nonnegative, Tonelli's theorem applies. Moreover, because distinct strips are disjoint, for every regular $t\in(\tau_0,T)$ we have
\[
\sum_S P_S(t)=\mathcal H^1(\{u=t\}\cap\Omega)\le P_V(t).
\]
Therefore
\begin{align}
\sum_S \operatorname{len}(\gamma_S)
&\le \frac{1}{2\pi}\sum_S\int_{a_S}^{b_S} P_S(t)\,dt \\
&= \frac{1}{2\pi}\int_{\tau_0}^{T}\sum_S \mathbf 1_{(a_S,b_S)}(t)P_S(t)\,dt \\
&= \frac{1}{2\pi}\int_{\tau_0}^{T}\sum_S P_S(t)\,dt \\
&\le \frac{1}{2\pi}\int_{\tau_0}^{T} P_V(t)\,dt \\
&\le \frac{1}{2\pi}\int_{2\alpha}^{\infty} P_V(t)\,dt.
\tag{11}
\end{align}

We now prove the uniqueness statements needed for the combinatorics.

\smallskip
\noindent\emph{Bottom strip.}
Set
\[
W_{\tau_0}:=V\cap\{u>\tau_0\}.
\]
We already know that $W_{\tau_0}$ is connected. Since there are no critical points on
\[
V\cap\{\tau_0\le u\le \mu_1-\delta_1\},
\]
Lemma~\ref{lem:regular-slab}(a), applied on $[\tau_0,\mu_1-\delta_1]$, shows that
\[
V^{\ge \mu_1-\delta_1}:=V\cap\{u\ge \mu_1-\delta_1\}
\]
is connected because $V^{\ge \tau_0}$ is connected. Choose a regular value
\[
\tau_1\in(\mu_1-\delta_1,\mu_1).
\]
Applying Lemma~\ref{lem:regular-slab}(a) on $[\mu_1-\delta_1,\tau_1]$ shows that
$V^{\ge \tau_1}$ is connected. Then Lemma~\ref{lem:regular-slab}(b), applied with
$a=\mu_1-\delta_1$ and $b=\tau_1$, implies that
\[
W_1:=V\cap\{u>\mu_1-\delta_1\}
\]
is connected.

By the component-inheritance paragraph above, both $W_{\tau_0}$ and $W_1$ are connected
components of the global regular superlevel sets $\{u>\tau_0\}$ and $\{u>\mu_1-\delta_1\}$,
respectively. Hence Lemma~\ref{lem:jordan-domain} applies to both of them, so each is a Jordan domain. Since
\[
\partial W_{\tau_0}\subset \{u=\tau_0\},
\qquad
\partial W_1\subset \{u=\mu_1-\delta_1\},
\]
these boundaries are disjoint. Because $W_1\subset W_{\tau_0}$, the auxiliary topological fact implies that
\[
\overline{W_1}\subset W_{\tau_0}
\]
and that
\[
W_{\tau_0}\setminus \overline{W_1}
\]
is an annulus. Since $\partial W_1\subset \{u=\mu_1-\delta_1\}$, we have
\[
A_0:=V\cap\{\tau_0<u<\mu_1-\delta_1\}=W_{\tau_0}\setminus \overline{W_1}.
\]
Therefore $A_0$ is a connected annulus. Since there are no critical points on $A_0$, Lemma~\ref{lem:regular-slab}(c) implies that every connected component of $A_0$ is a regular strip. Because $A_0$ is connected, it is a single regular strip. Hence there is exactly one strip whose lower cross-section lies in $\{u=\tau_0\}$; denote it by $S_0$.

\smallskip
\noindent\emph{Critical neighborhoods.}
\medskip
\noindent\emph{Claim.}
For the pairwise disjoint neighborhoods $N_r$ chosen above, each of the cross-sections
$C_r^{-}$, $C_{r,1}^{+}$, and $C_{r,2}^{+}$ lies on the boundary of exactly one connected component of
$\Omega$. More precisely, there exist unique connected components
\[
S_r^{-},\qquad S_{r,1}^{+},\qquad S_{r,2}^{+}
\]
of $\Omega$ such that $C_r^{-}$ is the upper cross-section of $S_r^{-}$ and
$C_{r,1}^{+}, C_{r,2}^{+}$ are the lower cross-sections of $S_{r,1}^{+}, S_{r,2}^{+}$, respectively.

\smallskip
\noindent\emph{Proof of claim.}
Fix $r$, and let $U_r$ and $N_r$ be the Morse chart and neighborhood chosen above. By properties \emph{(ii)} and \emph{(iii)}, the three cross-sections
\[
C_r^- \subset \{u=\mu_r-\delta_r\},
\qquad
C_{r,1}^+,\, C_{r,2}^+ \subset \{u=\mu_r+\delta_r\}
\]
appear exactly as in the standard saddle picture, with the remainder of $\partial N_r$
consisting of flow segments of $X$. By property \emph{(vi)}, there exist connected
open collar neighborhoods
\[
\mathcal C_r^-,
\quad
\mathcal C_{r,1}^+,
\quad
\mathcal C_{r,2}^+
\subset (U_r \setminus N_r)\cap \{\tau_0<u<T\}.
\]
Because the neighborhoods $U_r$ are pairwise disjoint and contain no other $N_k$, these collars are in fact contained in $\Omega$.

We first prove uniqueness for $C_r^-$. Suppose that two distinct connected components
$S_1,S_2$ of $\Omega$ both had $C_r^-$ in their boundary. Choose $q\in C_r^-$. By the
collar property, there exists $\rho>0$ such that
\[
B(q,\rho)\cap\Omega \subset \mathcal C_r^-.
\]
Since $q\in \partial S_1 \cap \partial S_2$, the ball $B(q,\rho)$ meets both $S_1$ and
$S_2$. But $\mathcal C_r^-$ is connected and contained in $\Omega$, so both $S_1$ and
$S_2$ must contain the connected set $\mathcal C_r^-$. Hence $S_1=S_2$, a contradiction.
Therefore $C_r^-$ lies on the boundary of exactly one connected component of $\Omega$;
denote it by $S_r^-$.

The same argument, applied to the collars $\mathcal C_{r,1}^+$ and $\mathcal C_{r,2}^+$,
shows that each of $C_{r,1}^+$ and $C_{r,2}^+$ lies on the boundary of exactly one
connected component of $\Omega$, denoted $S_{r,1}^+$ and $S_{r,2}^+$, respectively.

Finally, the local saddle model shows that $\mathcal C_r^-$ lies on the side from which
the flow enters the saddle, whereas $\mathcal C_{r,1}^+$ and $\mathcal C_{r,2}^+$ lie on
the two sides from which the flow exits. Hence $C_r^-$ is the upper cross-section of
$S_r^-$, and $C_{r,1}^+, C_{r,2}^+$ are the lower cross-sections of $S_{r,1}^+, S_{r,2}^+$,
respectively. This proves the claim.
\hfill$\triangle$
\smallskip
\noindent\emph{Top strips.}
Fix $j\in\{1,\dots,m\}$ and set
\[
c_j:=\mu_s+\delta_s.
\]
Let $W_j$ be the connected component of
\[
V\cap\{u>c_j\}
\]
that contains $K_j$, and set
\[
Y_j:=W_j\cap\{u\ge T\}.
\]
By the component-inheritance paragraph above, $W_j$ is also a connected component of the global regular superlevel set $\{u>c_j\}$. Hence Lemma~\ref{lem:jordan-domain} implies that $W_j$ is a Jordan domain.

We claim that $Y_j$ is connected.

For $z\in W_j$ with $u(z)<T$, set
\[
s_z:=u(z)\in(c_j,T)
\]
and
\[
M_{j,z}:=W_j\cap\{s_z\le u\le T\}.
\]
Because $s_z>c_j$ and $W_j$ is relatively closed in the open set $\{u>c_j\}$, the set
$M_{j,z}$ is closed in $\mathbb C$. It is also bounded because $\Lambda_f$ is bounded,
hence compact.

There are no critical points of $u$ on $M_{j,z}$, so the vector field
\[
X:=\frac{\nabla u}{|\nabla u|^2}
\]
is smooth on a neighborhood of $M_{j,z}$.

Let $\varphi_\tau(z)$ be the maximal trajectory of $X$ through $z$. As long as the
trajectory is defined,
\[
\frac{d}{d\tau}u(\varphi_\tau(z))=1,
\]
hence
\[
u(\varphi_\tau(z))=u(z)+\tau=s_z+\tau.
\]
Therefore, for every
\[
0\le \tau\le T-s_z,
\]
the point $\varphi_\tau(z)$ lies in $M_{j,z}$. Since $X$ is smooth on a neighborhood of the
compact set $M_{j,z}$, the continuation theorem for ODEs implies that the trajectory extends
throughout the whole interval $[0,T-s_z]$.

Define
\[
H_j:[0,1]\times W_j\to W_j,
\qquad
H_j(\theta,z):=
\begin{cases}
\varphi_{\theta(T-u(z))}(z),& u(z)<T,\\[1mm]
z,& u(z)\ge T.
\end{cases}
\]
To prove continuity, fix $s$ with $c_j<s<T$. On the compact slab
\[
W_j\cap\{s\le u\le T\},
\]
the vector field $X$ is smooth on a neighborhood, so its flow depends continuously on
initial data and time. Hence $H_j$ is continuous on
\[
[0,1]\times \bigl(W_j\cap\{u\ge s\}\bigr).
\]
Since $s\in(c_j,T)$ is arbitrary, these local continuity statements patch together to show
that $H_j$ is continuous on $[0,1]\times W_j$.

Thus $H_j$ is a strong deformation retraction of $W_j$ onto $Y_j$, and therefore $Y_j$
is connected.

Since $T$ is a regular value, the connected components of $V\cap\{u\ge T\}$ are precisely $K_1,\dots,K_m$. The set $Y_j\subset V\cap\{u\ge T\}$ is connected and contains $K_j$, hence $Y_j\subset K_j$. The reverse inclusion is immediate, so $Y_j=K_j$. In particular $Z_j$ is a Jordan domain and $K_j=\overline{Z_j}$. Since
\[
\partial W_j\subset \{u=c_j\},
\qquad
\partial Z_j=\partial K_j\subset \{u=T\},
\]
these boundaries are disjoint, and $Z_j\subset W_j$. Hence, by the auxiliary topological fact,
\[
K_j=\overline{Z_j}\subset W_j
\]
and
\[
W_j\setminus K_j=W_j\setminus \overline{Z_j}
\]
is an annulus. Because $\partial K_j\subset \{u=T\}$, we obtain
\[
A_j:=W_j\cap\{c_j<u<T\}=W_j\setminus K_j.
\]
Therefore $A_j$ is a connected annulus. Since there are no critical points on $A_j$, Lemma~\ref{lem:regular-slab}(c) implies that every connected component of $A_j$ is a regular strip. Because $A_j$ is connected, it is a single regular strip. Hence there is exactly one strip of $\Omega$ whose upper cross-section is $\partial K_j$; denote it by $S_j^{\mathrm{top}}$.

\medskip
\noindent\textbf{Exhaustion lemma for the strips.}
Every connected component $S$ of $\Omega$ has its lower cross-section in exactly one of
\[
\{u=\tau_0\}, \qquad C_{r,1}^+,\, C_{r,2}^+ \quad (1\le r\le s),
\]
and its upper cross-section in exactly one of
\[
C_r^- \quad (1\le r\le s), \qquad \partial K_j \quad (1\le j\le m).
\]
Conversely, the unique strip with lower cross-section in $\{u=\tau_0\}$ is $S_0$;
the unique strips with lower cross-sections $C_{r,1}^+,C_{r,2}^+$ are
$S_{r,1}^+,S_{r,2}^+$; the unique strips with upper cross-sections $C_r^-$ are
$S_r^-$; and the unique strips with upper cross-sections $\partial K_j$ are
$S_j^{\mathrm{top}}$.

\smallskip
\noindent\emph{Proof.}
Let $S$ be a connected component of $\Omega$. By Lemma~\ref{lem:regular-strip}, $S$ is a regular strip,
so it has exactly one lower cross-section and exactly one upper cross-section, both lying
on boundary components of $\Omega$ contained in level sets of $u$.
Now $\partial\Omega$ consists of:
\begin{itemize}
\item the base level $\{u=\tau_0\}$,
\item the top boundaries $\partial K_j\subset\{u=T\}$,
\item the cross-sections $C_r^-\subset\{u=\mu_r-\delta_r\}$ and
      $C_{r,1}^+,C_{r,2}^+\subset\{u=\mu_r+\delta_r\}$,
\item and side-boundary flow segments on the sets $\partial N_r$.
\end{itemize}
By definition, a cross-section of a regular strip cannot lie on a side-boundary flow
segment. Since the lower level of $S$ is strictly smaller than its upper level, the lower
cross-section cannot lie on any $\partial K_j$ or $C_r^-$, and the upper cross-section
cannot lie on $\{u=\tau_0\}$ or any $C_{r,i}^+$. This proves the classification.

The converse uniqueness statements are exactly the uniqueness assertions for
bottom strips, critical neighborhoods, and top strips proved above.
\hfill$\triangle$

We now define a finite abstract graph $T$. Its vertices are
\[
V(T)=\{v_0\}\cup\{w_1,\dots,w_s\}\cup\{\ell_1,\dots,\ell_m\},
\]
where $v_0$ is the root, $w_r$ corresponds to the saddle $p_r$, and $\ell_j$
corresponds to the zero $a_j$.

For each strip $S\subset\Omega$, define an oriented edge $e_S$ from the vertex attached
to its lower cross-section to the vertex attached to its upper cross-section:
\begin{itemize}
\item if the lower cross-section of $S$ lies in $\{u=\tau_0\}$, then the lower endpoint
      of $e_S$ is $v_0$;
\item if the lower cross-section of $S$ is one of the arcs $C_{r,i}^+$, then the lower
      endpoint of $e_S$ is $w_r$;
\item if the upper cross-section of $S$ is $C_r^-$, then the upper endpoint of $e_S$ is
      $w_r$;
\item if the upper cross-section of $S$ is $\partial K_j$, then the upper endpoint of
      $e_S$ is $\ell_j$.
\end{itemize}
By the exhaustion lemma above, every strip determines exactly one such edge.

Every non-root vertex has exactly one incoming edge:
$w_r$ receives the edge $e_{S_r^-}$, and $\ell_j$ receives the edge
$e_{S_j^{\mathrm{top}}}$.
Likewise each $w_r$ has exactly two outgoing edges, namely
$e_{S_{r,1}^+}$ and $e_{S_{r,2}^+}$.

Assign heights by
\[
h(v_0)=\tau_0,\qquad h(w_r)=\mu_r,\qquad h(\ell_j)=T+1.
\]
Every oriented edge goes from a vertex of smaller height to a vertex of larger height.
Hence there is no directed cycle. Moreover, iterating the unique predecessor map from
any non-root vertex strictly decreases height, so after finitely many steps one reaches
$v_0$. Thus every vertex is joined to $v_0$ by an undirected path, and the underlying
graph is connected.

We now show that the underlying graph is acyclic. By the construction above, every
non-root vertex has exactly one incoming edge, while $v_0$ has none. Therefore
\[
|E(T)| = |V(T)| - 1.
\]
Since the underlying graph is finite and connected, this implies that it is a tree. Hence
the underlying graph of $T$ is acyclic, and therefore $T$ is a finite rooted tree.

We now realize \(T\) geometrically. The neighborhoods \(N_r\) and the value \(T\) have already been fixed, and hence so have \(\Omega\), its strips \(S\), and the trajectory segments \(\gamma_S\).

For each \(r\), let
\[
x_{r,-}:=\gamma_{S_r^-}\cap C_r^-,
\qquad
x_{r,i}:=\gamma_{S_{r,i}^+}\cap C_{r,i}^+
\quad (i=1,2).
\]
These points are well defined, because \(S_r^-\) has upper cross-section \(C_r^-\) and \(S_{r,i}^+\) has lower cross-section \(C_{r,i}^+\). They are pairwise distinct, since the three cross-sections
\[
C_r^-,\ C_{r,1}^+,\ C_{r,2}^+
\]
are disjoint subsets of \(\partial N_r\). Set
\[
c_r:=p_r.
\]
By property \((v)\) for the chosen neighborhood $N_r$, there exists an embedded \(Y\)-shaped tree
\[
T_r:=T_r(x_{r,-},x_{r,1},x_{r,2})\subset N_r
\]
with branch point \(c_r\) and terminal vertices \(x_{r,-},x_{r,1},x_{r,2}\), and
\[
\operatorname{len}(T_r)\le \eta_r.
\]
Therefore
\[
\sum_{r=1}^s \operatorname{len}(T_r)<\frac{\varepsilon}{2}. \tag{12}
\]

Let \(x_0\) be the lower endpoint of \(\gamma_{S_0}\) on the level \(\{u=\tau_0\}\). For each \(j\), let \(b_j\) be the upper endpoint of \(\gamma_{S_j^{\mathrm{top}}}\) on \(\partial K_j\). By the earlier choice of $B(a_j,r_j)$, $\psi_j$, and $T$, we have
\[
Z_j=\psi_j^{-1}\bigl(D(0,e^{-T/m_j})\bigr),
\qquad
K_j=\psi_j^{-1}\bigl(\overline{D(0,e^{-T/m_j})}\bigr)\subset B(a_j,r_j).
\]
If $w_j:=\psi_j(b_j)$, then $|w_j|=e^{-T/m_j}$, so the radial segment $[0,w_j]$
is contained in $\overline{D(0,e^{-T/m_j})}$ and its relative interior is contained in
$D(0,e^{-T/m_j})$. Therefore its preimage under $\psi_j^{-1}$ is an embedded arc
\[
\sigma_j\subset K_j,\qquad \sigma_j^\circ\subset Z_j,
\]
joining $a_j$ to $b_j$. By the choice of $C_j$,
\[
\operatorname{len}(\sigma_j)\le C_j e^{-T/m_j}.
\]
Therefore, by the choice of $T$,
\[
\sum_{j=1}^m \operatorname{len}(\sigma_j)<\frac{\varepsilon}{2}. \tag{13}
\]

\indent For each strip $S$, define an embedded arc $A_S$ as follows:
\begin{itemize}
\item if $e_S$ joins $v_0$ to $w_r$, let $A_S$ be the union of $\gamma_S$ and the unique arm of $T_r$ joining the upper
      endpoint of $\gamma_S$ to $c_r$;
\item if $e_S$ joins $w_r$ to $w_{r'}$, let $A_S$ be the union of the unique arm of $T_r$ from $c_r$ to the lower endpoint
      of $\gamma_S$, the segment $\gamma_S$, and the unique arm of $T_{r'}$ from the upper endpoint of $\gamma_S$ to $c_{r'}$;
\item if $e_S$ joins $w_r$ to $\ell_j$, let $A_S$ be the union of the unique arm of $T_r$ from $c_r$ to the lower endpoint
      of $\gamma_S$, the segment $\gamma_S$, and the arc $\sigma_j$ from the upper endpoint of $\gamma_S$ to $a_j$.
\end{itemize}
The geometric vertices are
\[
 p_{v_0}=x_0,\qquad p_{w_r}=c_r,\qquad p_{\ell_j}=a_j.
\]
We verify the hypotheses of Lemma~\ref{lem:geometric-tree-realization} for the family of arcs $\{A_S\}$. First, if $S\neq S'$, then
\[
\gamma_S\cap\gamma_{S'}=\varnothing,
\]
because distinct strips are disjoint subsets of $\Omega$. Second, for each $r$, the tree $T_r\subset N_r$ meets the rest of the construction only at its terminal points $x_{r,-},x_{r,1},x_{r,2}$: indeed the neighborhoods $N_r$ are pairwise disjoint, each $N_r$ is disjoint from every $Z_j$, and a strip trajectory can meet $\partial N_r$ only at the cross-section belonging to that strip. Third, for each $j$, the arc $\sigma_j\subset K_j$ has interior contained in $Z_j$, and it meets the rest of the construction only at its endpoint $b_j\in\partial K_j$. Indeed, the sets $Z_1,\dots,Z_m$ are pairwise disjoint and disjoint from all $N_r$, so the arcs $\sigma_j$ are pairwise disjoint and are disjoint from every $T_r$. Moreover, every strip trajectory $\gamma_S$ is contained in
\[
\Omega\subset\{\tau_0<u<T\},
\]
whereas
\[
\sigma_j^\circ\subset Z_j\subset\{u>T\}.
\]
Hence $\gamma_S\cap \sigma_j^\circ=\varnothing$, and the only possible intersection between $\sigma_j$ and the strip part of the construction is the point $b_j$, where $\sigma_j$ attaches to $\gamma_{S_j^{\mathrm{top}}}$.

It follows that the interior of each $A_S$ contains no geometric vertex, the interiors of distinct arcs $A_S$ and $A_{S'}$ are disjoint, and two distinct arcs can meet only at a common geometric endpoint corresponding to a common endpoint of the associated edges of $T$. Hence Lemma~\ref{lem:geometric-tree-realization} applies:
\[
 H:=\bigcup_S A_S
\]
is homeomorphic to $T$. Since $T$ is a tree, $H$ is a finite embedded tree in $V$ containing all zeros
$a_1,\dots,a_m$.

\indent Finally,
\[
\operatorname{len}(H)\le \sum_S \operatorname{len}(\gamma_S)+\sum_{r=1}^s \operatorname{len}(T_r)+\sum_{j=1}^m \operatorname{len}(\sigma_j).
\]
Using (11), (12), and (13), we obtain
\[
\operatorname{len}(H)\le \frac{1}{2\pi}\int_{2\alpha}^{\infty} P_V(t)\,dt+\varepsilon.
\]
Setting $G_\varepsilon:=H$ proves (10). \end{proof}

\section{An explicit perturbation lemma}

We now replace the abstract transversality step by a concrete two-parameter perturbation. Fix a
squarefree monic polynomial $f$, a component $U\subset \Lambda_f$ containing exactly $m\ge 2$ zeros of $f$, and the
number $\alpha$ introduced above. Define the compact collars
\[
K:=\{z\in U:u(z)\ge \alpha/2\},\qquad L:=\{z\in U:u(z)\ge 3\alpha/2\}.
\]
Then $K$ and $L$ are compact and connected by Lemma~\ref{lem:superlevel-connected}.

\indent Set
\[
 c:=3\alpha/2.
\]
Since $u$ has no critical values in $[0,4\alpha)$, the value $c$ is regular. Let
\[
W_1,\dots,W_r
\]
be the connected components of $\{z\in U:u(z)>c\}$. We claim that the closures $\overline{W_i}$ are pairwise
disjoint. Indeed, if $p\in \overline{W_i}\cap\overline{W_j}$, then $u(p)=c$. Because $c>0$, the point $p$ lies in the interior of $U$.
Since $c$ is a regular value, the implicit function theorem gives local coordinates near $p$ in which
\[
 u-c=x.
\]
Hence locally $\{u>c\}$ is a single open half-disk, so only one connected component of $\{z\in U:u(z)>c\}$ can
accumulate at $p$. Thus the sets $\overline{W_i}$ are pairwise disjoint. Every point of
\[
 L=\{z\in U:u(z)\ge c\}
\]
lies in $\bigcup_i \overline{W_i}$: if $u(p)>c$, then $p\in\{u>c\}=\bigcup_i W_i\subset
\bigcup_i \overline{W_i}$; if $u(p)=c$, then the same local normal form shows that points with
$u>c$ lie arbitrarily close to $p$, so again $p\in\bigcup_i \overline{W_i}$. Therefore
\[
L=\bigcup_{i=1}^r \overline{W_i}.
\]
Since $L$ is connected by Lemma~\ref{lem:superlevel-connected} and the $\overline{W_i}$ are pairwise disjoint closed subsets of $L$, it follows
that $r=1$. Let $W:=W_1$. Because $W\subset U\subset\{u>0\}$ and $c>0$, the set $W$ is also a connected
component of the global superlevel set $\{u>c\}$. Hence Lemma~\ref{lem:jordan-domain} implies that $W$ is a Jordan domain and
\[
L=\overline W.
\]
An entirely analogous argument with
\[
 d:=\alpha/2
\]
shows that $K$ is a compact Jordan domain. Indeed, $d$ is also a regular value because $u$ has no critical
values in $[0,4\alpha)$. If $Y_1,\dots,Y_s$ are the connected components of $\{u>d\}$, then exactly as above their
closures are pairwise disjoint and
\[
K=\{u\ge d\}=\bigcup_{i=1}^s \overline{Y_i}.
\]
Since $K$ is connected by Lemma~\ref{lem:superlevel-connected}, we must have $s=1$. Let $Y:=Y_1$. Because $d>0$ and $Y\subset U\subset\{u>0\}$,
this set $Y$ is a connected component of the global superlevel set $\{u>d\}$. Hence Lemma~\ref{lem:jordan-domain} implies that $Y$
is a Jordan domain and
\[
K=\overline Y.
\]
Choose pairwise disjoint closed disks
\[
D_1,\dots,D_m\Subset \{u>5\alpha/2\}\cap W
\]
centered at the zeros of $f$ in $U$, so small that each $D_j$ contains exactly one
zero and
\[
\delta_D:=\min_{z\in \bigcup_{j=1}^m D_j}|f'(z)|>0.
\]
In particular,
\[
B:=\{z\in U:\alpha/2\le u(z)\le 5\alpha/2\}
\]
is disjoint from $\bigcup_{j=1}^m D_j$. Define
\[
L_0:=L\setminus \bigcup_{j=1}^m \operatorname{int}(D_j).
\]

Since $W$ is a Jordan domain, $\overline W$ is homeomorphic to the closed unit disk
$\overline{\mathbb D}$ by the Jordan--Schoenflies theorem. Under such a homeomorphism,
the sets $D_j$ correspond to pairwise disjoint compact topological disks
$E_j\subset \mathbb D$.

\medskip
\noindent\emph{Auxiliary fact.}
Let $\Omega$ be a closed topological disk, and let
$E_1,\dots,E_m\subset \operatorname{int}(\Omega)$ be pairwise disjoint compact
topological disks. Then
\[
\Omega\setminus \bigcup_{j=1}^m \operatorname{int}(E_j)
\]
is path connected.

\smallskip
\noindent\emph{Proof of the auxiliary fact.}
Set
\[
X_r:=\Omega\setminus \bigcup_{j=1}^r \operatorname{int}(E_j)
\qquad (0\le r\le m),
\]
so $X_0=\Omega$. We first prove by induction on $r$ that each $X_r$ is connected.
This is clear for $r=0$. Assume that $X_{r-1}$ is connected, and consider
\[
X_r=X_{r-1}\setminus \operatorname{int}(E_r).
\]
Let $U$ be a connected component of $X_r\setminus \partial E_r=X_{r-1}\setminus E_r$.
Then $U$ is a nonempty proper open subset of the connected space $X_{r-1}$, so its
boundary in $X_{r-1}$ is nonempty. Every such boundary point must lie in $E_r\cap X_r=\partial E_r$.
Hence the closure of each connected component of $X_r\setminus \partial E_r$ in $X_r$
meets $\partial E_r$.

Now suppose that $X_r=A\cup B$ is a separation. Since $\partial E_r$ is connected, it must lie
entirely in one side, say $\partial E_r\subset A$. Each connected component $U$ of
$X_r\setminus \partial E_r$ is contained either in $A$ or in $B$; but if $U\subset B$, then,
because $B$ is closed in $X_r$, the closure of $U$ in $X_r$ would also be contained in $B$,
contradicting $\overline U\cap \partial E_r\neq\varnothing$. Thus every such component lies in $A$,
and therefore $B=\varnothing$. So $X_r$ is connected, completing the induction.

Finally, each point of $X_m$ has a neighborhood in $X_m$ homeomorphic either to an open disk
or to a closed half-disk. For interior points this is immediate. If $x$ lies on $\partial\Omega$
or on some $\partial E_j$, then, because the finitely many compact Jordan curves
$\partial\Omega,\partial E_1,\dots,\partial E_m$ are pairwise disjoint, there is a small neighborhood
of $x$ meeting no boundary curve other than the one through $x$. The Jordan--Schoenflies theorem
then straightens that boundary curve locally to a line, so the corresponding neighborhood of $x$ in
$X_m$ is homeomorphic to a closed half-disk. Thus $X_m$ is locally path connected. Since connected
and locally path connected spaces are path connected, $X_m$ is path connected. This proves the auxiliary fact.
\medskip

Applying this with $\Omega=\overline{\mathbb D}$, we obtain that
\[
\overline{\mathbb D}\setminus \bigcup_{j=1}^m \operatorname{int}(E_j)
\]
is path connected. Pulling paths back under the Schoenflies homeomorphism shows that
\[
L_0
=
L\setminus \bigcup_{j=1}^m \operatorname{int}(D_j)
\]
is path connected, hence connected.

For $\lambda,\beta\in\mathbb{C}$ define
\[
 g_{\lambda,\beta}(z):=f(z)+\lambda z+\beta,\qquad u_{\lambda,\beta}:=-\log|g_{\lambda,\beta}|.
\]

\medskip
\begin{lemma}[Explicit perturbation]\label{lem:explicit-perturbation}
\emph{There exist arbitrarily small parameters $(\lambda,\beta)\in\mathbb{C}^2$ such that,
writing $g:=g_{\lambda,\beta}$ and $u_g:=-\log|g|$, the following hold:}

\noindent (1) \emph{each $D_j$ contains exactly one zero of $g$, this zero is simple, and $g$ has no zeros on
$K\setminus \bigcup_j D_j$;}

\noindent (2) \emph{$u_g$ has no critical values in $[\alpha,2\alpha]$ on $K$;}

\noindent (3) \emph{every critical point of $u_g$ in $K$ with value $>2\alpha$ is nondegenerate;}

\noindent (4) \emph{all critical values of $u_g$ in $K$ larger than $2\alpha$ are pairwise distinct.}
\end{lemma}

\begin{proof} We proceed in five steps.

Step 1: keep the zeros in the prescribed disks. Since each \(D_j\) contains exactly one simple zero of
\(f\) and \(f\) has no zeros on
\[
C:=K\setminus \bigcup_{j=1}^m \operatorname{int}(D_j),
\]
we have
\[
\delta_C:=\min_{z\in C}|f(z)|>0.
\]
Choose \(\eta_0>0\) so small that \(\eta_0<\delta_C/2\). If
\[
\sup_{z\in K}|g(z)-f(z)|<\eta_0,
\]
then \(g\) has no zeros on \(C\), since
\[
|g(z)|\ge |f(z)|-|g(z)-f(z)|>\delta_C/2
\qquad (z\in C).
\]
On the other hand, for each \(j\), since \(f\) has exactly one zero in \(D_j\) and no zero on
\(\partial D_j\), Rouché's theorem on \(\partial D_j\) shows that \(g\) also has exactly one zero in
\(D_j\), counted with multiplicity, provided the same uniform bound holds. Thus each \(D_j\) still
contains exactly one zero of \(g\), and \(g\) has no zeros on
\[
K\setminus \bigcup_{j=1}^m D_j.
\]

Since
\[
\sup_{z\in K}|g_{\lambda,\beta}(z)-f(z)|
\le
|\lambda|\max_{z\in K}|z|+|\beta|,
\]
this zero configuration is preserved for all sufficiently small \((\lambda,\beta)\).

Moreover, if
\[
|\lambda|<\delta_D/2,
\]
then for every \(z\in \bigcup_{j=1}^m D_j\),
\[
|g'(z)|=|f'(z)+\lambda|\ge |f'(z)|-|\lambda|>\delta_D/2,
\]
so \(g'\) has no zeros on any \(D_j\). Combined with the preceding paragraph, this shows that the
unique zero of \(g\) in each \(D_j\) is simple.

\noindent\textit{Step 2: make the critical points simple.} 


On the zero-free set $\{g\neq 0\}$, the function $u_g=-\log|g|$ is smooth, and its critical points are exactly the zeros of $g'$, because
\[
\nabla(-\log|g|)=0 \iff g'=0
\qquad\text{on }\{g\neq 0\}.
\]
But
\[
g'(z)=f'(z)+\lambda.
\]
A zero $\zeta$ of $f'+\lambda$ is multiple if and only if
\[
f'(\zeta)+\lambda=0
\quad\text{and}\quad
f''(\zeta)=0,
\]
equivalently,
\[
\lambda=-f'(\zeta)
\quad\text{with}\quad
f''(\zeta)=0.
\]
Therefore the bad set
\[
\Sigma_1:=\{-f'(\zeta):f''(\zeta)=0\}
\]
is finite, and for every $\lambda\notin\Sigma_1$ all zeros of $f'+\lambda$ are simple.

We now show that simplicity of the zero of $g'$ implies nondegeneracy of the corresponding critical point of $u_g$. Let $c$ satisfy $g(c)\neq 0$, $g'(c)=0$, and $g''(c)\neq 0$. Choose a simply connected neighborhood of $c$ on which $g$ has no zeros, and let $F=\log g$ be a holomorphic branch there. Then
\[
u_g=-\Re F.
\]
Since
\[
F'(z)=\frac{g'(z)}{g(z)},
\]
we have $F'(c)=0$, and
\[
F''(c)=\frac{g''(c)}{g(c)}\neq 0.
\]
Hence
\[
F(z)=F(c)+\frac{F''(c)}{2}(z-c)^2+O(|z-c|^3),
\]
so
\[
u_g(z)
=
u_g(c)-\Re\!\left(\frac{F''(c)}{2}(z-c)^2\right)+O(|z-c|^3).
\]
The quadratic form $\Re\!\big(A(z-c)^2\big)$ with $A\neq 0$ is nondegenerate and has signature $(1,1)$. Therefore the Hessian of $u_g$ at $c$ is nondegenerate. Thus every critical point of $u_g$ on $\{g\neq 0\}$ corresponding to a simple zero of $g'$
is a Morse saddle.

Consequently, for every $\lambda\notin\Sigma_1$, every critical point of $u_g$ with finite critical value is nondegenerate.



\noindent\textit{Step 3: make the complex critical values distinct.} Set
\[
H_\lambda(z):=f(z)+\lambda z.
\]
Its critical points are the zeros of
\[
H_\lambda'(z)=f'(z)+\lambda.
\]
The critical values of $H_\lambda$ are contained among the roots of the resultant polynomial
\[
R_\lambda(w):=\operatorname{Res}_z\bigl(f'(z)+\lambda,\,f(z)+\lambda z-w\bigr).
\]
If two distinct critical points of $H_\lambda$ have the same critical value, then
$R_\lambda$ has a repeated root. Hence the discriminant
\[
\Delta(\lambda):=\operatorname{Disc}_w R_\lambda(w)
\]
vanishes. Therefore every $\lambda$ for which the critical values of $H_\lambda$ are not
pairwise distinct belongs to the algebraic set
\[
\Sigma_2:=\{\lambda:\Delta(\lambda)=0\}.
\]
To prove that such bad parameters do not fill the plane, it suffices to exhibit one value of
$\lambda$ for which the critical values of $H_\lambda$ are pairwise distinct. The asymptotic
argument below does exactly this: for all sufficiently large $|\lambda|$, the critical points $c_k(\lambda)$
satisfy
\[
c_k(\lambda)=\rho\omega_k+O(1),
\]
and the corresponding critical values satisfy
\[
H_\lambda(c_k(\lambda))=-(n-1)\rho^n\omega_k+O(|\rho|^{n-1}),
\]
so they are pairwise distinct. We now carry out that asymptotic argument in detail.
Write
\[
 f'(z)=nz^{n-1}+Q(z),\qquad \deg Q\le n-2.
\]
Fix $R>0$, to be chosen later. Let $\lambda\neq 0$, choose $\rho\in\mathbb{C}$ such that
\[
 \rho^{n-1}=-\lambda/n,
\]
and let $\omega_1,\dots,\omega_{n-1}$ denote the $(n-1)$st roots of unity. For each $k$ set
\[
 c_k^0:=\rho\omega_k,\qquad C_k:=\{z:|z-c_k^0|=R\}.
\]
Then $c_1^0,\dots,c_{n-1}^0$ are precisely the zeros of $nz^{n-1}+\lambda$.

\indent Let $z\in C_k$, so that $z=c_k^0+\xi$ with $|\xi|=R$. Since
\[
 nz^{n-1}+\lambda=nz^{n-1}-(c_k^0)^{n-1}=n(z-c_k^0)\sum_{m=0}^{n-2} z^{n-2-m}(c_k^0)^m,
\]
we have
\[
 |nz^{n-1}+\lambda|=nR\left|\sum_{m=0}^{n-2} z^{n-2-m}(c_k^0)^m\right|.
\]
Factoring out $(c_k^0)^{n-2}$ gives
\[
\sum_{m=0}^{n-2} z^{n-2-m}(c_k^0)^m=(c_k^0)^{n-2}\sum_{m=0}^{n-2}\left(1+\frac{\xi}{c_k^0}\right)^{n-2-m}.
\]
Because $R$ is fixed and $|c_k^0|=|\rho|\to\infty$ as $|\lambda|\to\infty$, the second factor converges uniformly on all the
circles $C_k$ to $n-1$. Therefore there exists $\rho_1>0$ such that, for all $|\rho|\ge \rho_1$ and all $z\in C_k$,
\[
\left|\sum_{m=0}^{n-2} z^{n-2-m}(c_k^0)^m\right|\ge \frac{n-1}{2}|\rho|^{n-2}.
\]
Hence
\begin{equation}
|nz^{n-1}+\lambda|\ge \frac{n(n-1)}{2}R|\rho|^{n-2}\qquad (z\in C_k,\ |\rho|\ge \rho_1).
\tag{14}
\end{equation}
Now write
\[
Q(z)=\sum_{j=0}^{n-2} a_j z^j
\]
and set
\[
 C_0:=\sum_{j=0}^{n-2}|a_j|\left(\frac32\right)^j.
\]
If $|\rho|\ge 2R$ and $z\in C_k$, then $|z|\le |\rho|+R\le \tfrac32|\rho|$, so
\begin{equation}
|Q(z)|\le \sum_{j=0}^{n-2}|a_j||z|^j\le \sum_{j=0}^{n-2}|a_j|\left(\frac32\right)^j |\rho|^j\le C_0 |\rho|^{n-2}.
\tag{15}
\end{equation}
Choose $R>0$ so large that
\[
\frac{n(n-1)}{2}R>C_0.
\]
Then, by (14) and (15), there exists $\rho_0\ge \max\{\rho_1,2R\}$ such that for all $|\rho|\ge \rho_0$ and all $z\in C_k$,
\[
 |Q(z)|<|nz^{n-1}+\lambda|.
\]
Therefore Rouch\'e's theorem implies that, for every such $\rho$, the polynomial
\[
 f'(z)+\lambda=nz^{n-1}+\lambda+Q(z)
\]
has exactly one zero inside each disk bounded by $C_k$. Since $f'(z)+\lambda$ has degree $n-1$, these are all
of its zeros. Denote them by $c_1(\lambda),\dots,c_{n-1}(\lambda)$. In particular,
\[
 c_k(\lambda)=\rho\omega_k+O(1)
\]
uniformly in $k$ as $|\lambda|\to\infty$.

\indent Now write
\[
 f(z)=z^n+b_{n-1}z^{n-1}+\cdots+b_0.
\]
Since $f'(c_k(\lambda))+\lambda=0$, we have
\[
H_\lambda(c_k(\lambda))=f(c_k(\lambda))+\lambda c_k(\lambda)=f(c_k(\lambda))-c_k(\lambda)f'(c_k(\lambda)).
\]
But
\[
 f(z)-zf'(z)=-(n-1)z^n+\widetilde Q(z),
\]
where $\deg \widetilde Q\le n-1$. Hence
\[
H_\lambda(c_k(\lambda))=-(n-1)c_k(\lambda)^n+\widetilde Q(c_k(\lambda)).
\]
Since $c_k(\lambda)=\rho\omega_k+O(1)$, it follows that
\[
 c_k(\lambda)^n=\rho^n\omega_k^n+O(|\rho|^{n-1})=\rho^n\omega_k+O(|\rho|^{n-1}),
\]
because $\omega_k^{n-1}=1$ implies $\omega_k^n=\omega_k$. Also
\[
 \widetilde Q(c_k(\lambda))=O(|\rho|^{n-1}).
\]
Therefore
\[
 H_\lambda(c_k(\lambda))=-(n-1)\rho^n\omega_k+O(|\rho|^{n-1}).
\]
If $k\neq \ell$, then
\[
 H_\lambda(c_k(\lambda))-H_\lambda(c_\ell(\lambda))=-(n-1)\rho^n(\omega_k-\omega_\ell)+O(|\rho|^{n-1}).
\]
Since $\omega_k\neq \omega_\ell$, the leading term has modulus bounded below by $c|\rho|^n$ for some constant $c>0$
depending only on $n$, whereas the error term is $O(|\rho|^{n-1})$. Consequently, for all sufficiently large
$|\lambda|$, the critical values $H_\lambda(c_1(\lambda)),\dots,H_\lambda(c_{n-1}(\lambda))$ are pairwise distinct.

\indent Thus $\Delta(\lambda)\neq 0$ for all sufficiently large $|\lambda|$, so $\Delta\not\equiv 0$. Since $\Delta$ is a polynomial, the set $\Sigma_2$
is finite. Choose a small $\lambda\notin \Sigma_1\cup\Sigma_2$. Then $H_\lambda$ has simple critical points and pairwise distinct
complex critical values.

\noindent\textit{Step 4: keep all critical values in $K$ above $2\alpha$.} Let
\[
 Z_K:=\{\zeta\in K:f'(\zeta)=0\}.
\]
Since $f'$ has no zeros on $\partial K$, the set $Z_K$ is finite and contained in $\operatorname{int}(K)$.
Set
\[
M:=\max_{z\in K}|z|.
\]
If $Z_K=\varnothing$, then $f'(z)\neq 0$ for every $z\in K$. Since $K$ is compact and
$|f'|$ is continuous, the minimum
\[
\delta_K:=\min_{z\in K}|f'(z)|
\]
is strictly positive. Choose $\lambda$ so small that
\[
|\lambda|<\frac{\delta_K}{2}.
\]
Then for every $z\in K$,
\[
|g'(z)|=|f'(z)+\lambda|\ge |f'(z)|-|\lambda|\ge \frac{\delta_K}{2}>0.
\]
Hence $g'$ has no zeros on $K$. In particular, $u_g$ has no critical points in $K$,
so there are no critical values of $u_g$ arising from critical points in $K$. Moreover, any
zero of $g$ lying in one of the root-disks $D_j$ is a logarithmic singularity of $u_g$, not a
smooth critical point, and by Step~1 it satisfies $u_g>5\alpha/2$. Therefore such zeros do not
contribute any level in $[\alpha,2\alpha]$. In particular, no critical value of $u_g$ on $K$
lies in $[\alpha,2\alpha]$. In this case set
\[
\beta_*:=1.
\]

If $Z_K\neq\varnothing$, then, because $u=-\log|f|$ has no critical values in $[0,4\alpha)$, every $\zeta\in Z_K$ satisfies
\[
 |f(\zeta)|\le e^{-4\alpha}.
\]
Set
\[
M_*:=\max_{\zeta\in Z_K}|f(\zeta)|,\qquad \eta_*:=e^{-2\alpha}-M_*>0.
\]

Choose pairwise disjoint closed disks $E_\zeta\Subset \operatorname{int}(K)$, one around each
$\zeta\in Z_K$, also disjoint from $\bigcup_{j=1}^m D_j$, so small that
\[
|f(z)|\le M_*+\frac{\eta_*}{4}=e^{-2\alpha}-\frac{3\eta_*}{4}
\qquad
\left(z\in E:=\bigcup_{\zeta\in Z_K}E_\zeta\right),
\]
and such that $f'$ has no zeros on $K\setminus E$. Set
\[
\delta_E:=\min_{z\in K\setminus E}|f'(z)|>0,
\qquad
\delta_*:=\min_{\zeta\in Z_K}\min_{z\in \partial E_\zeta}|f'(z)|>0.
\]

Shrink $\lambda$ further so that
\[
|\lambda|<\min\left\{\frac{\delta_E}{2},\,\frac{\delta_*}{2},\,\frac{\eta_*}{4M}\right\}.
\]
Then for every $z\in K\setminus E$,
\[
|f'(z)+\lambda|\ge |f'(z)|-|\lambda|>\frac{\delta_E}{2}>0,
\]
so $f'+\lambda$ has no zeros on $K\setminus E$.

Moreover, for each $\zeta\in Z_K$ and every $z\in \partial E_\zeta$,
\[
|\lambda|<|f'(z)|.
\]
Hence, by Rouch\'e's theorem on $\partial E_\zeta$, the functions $f'$ and $f'+\lambda$
have the same number of zeros in $E_\zeta$, counted with multiplicity. In particular,
every zero of $f'+\lambda$ lying in $K$ belongs to $E$. Since the critical points of
$H_\lambda(z):=f(z)+\lambda z$ are precisely the zeros of
\[
H_\lambda'(z)=f'(z)+\lambda,
\]
it follows that every critical point $c$ of $H_\lambda$ in $K$ lies in $E$. Therefore
\[
|f(c)|\le e^{-2\alpha}-\frac{3\eta_*}{4}.
\]
Hence
\[
|H_\lambda(c)|
\le |f(c)|+|\lambda|\,|c|
\le |f(c)|+|\lambda|M
<
e^{-2\alpha}-\frac{\eta_*}{2}.
\]

Now choose $\beta$ so small that
\[
|\beta|<\frac{\eta_*}{2}.
\]
Then for every critical point $c$ of $H_\lambda$ in $K$,
\[
|g(c)|=|H_\lambda(c)+\beta|<e^{-2\alpha}.
\]
Thus every zero $c$ of $g'$ in $K$ satisfies $|g(c)|<e^{-2\alpha}$.
After proving below that no such $c$ is a zero of $g$, it will follow that every
critical value of $u_g$ arising from a critical point in $K$ is $>2\alpha$.
Consequently, once that nonvanishing statement is established, no critical point
of $u_g$ in $K$ can have critical value in $[\alpha,2\alpha]$.
In this case set
\[
\beta_*:=\frac{\eta_*}{2}.
\]

\noindent\textit{Step 5: make the moduli distinct.} We first dispose of the case $Z_K=\varnothing$.
In that case, by the first part of Step~4 we chose $\lambda$ so small that
\[
|g'(z)| = |f'(z)+\lambda| \ge \delta_K/2 > 0
\qquad (z\in K),
\]
so $g'$ has no zeros on $K$. Hence $g$ has no critical points on $K$, and therefore
$u_g$ has no critical values arising from critical points in $K$. In particular,
conclusion~(4) is vacuous in this case. We may therefore choose any $\beta$ with
\[
0<|\beta|<\min\{\beta_1,\beta_*\},
\]
where $\beta_1$ is as in Step~1 and $\beta_*=1$ is the choice made in Step~4, and
the lemma is proved in this case.

Assume from now on that $Z_K\neq\varnothing$. We now verify the deferred claim
that no critical point of $g$ in $K$ is a zero of $g$. Indeed, if $c\in K$ satisfies
\[
g'(c)=0,
\]
then
\[
f'(c)+\lambda = 0,
\]
so $c$ is a critical point of
\[
H_\lambda(z) := f(z)+\lambda z.
\]
By Step~4, every critical point of $H_\lambda$ in $K$ lies in
\[
E := \bigcup_{\zeta\in Z_K} E_\zeta \subset K \setminus \bigcup_{j=1}^m D_j.
\]
By Step~1, $g$ has no zeros on
\[
K \setminus \bigcup_{j=1}^m D_j.
\]
Hence
\[
g(c)\neq 0
\]
for every critical point $c$ of $g$ in $K$. Therefore each such $c$ is a genuine smooth
critical point of $u_g$, and its critical value is
\[
u_g(c) = -\log|g(c)| = -\log|H_\lambda(c)+\beta|.
\]

Let $a_1,\dots,a_N$ be the pairwise distinct complex critical
values of $H_\lambda$ arising from critical points in $K$. After adding $\beta$, they become $a_1+\beta,\dots,a_N+\beta$.
Their moduli fail to be distinct only if, for some $i\neq j$,
\[
 |a_i+\beta|=|a_j+\beta|.
\]
Equivalently,
\[
 |a_i+\beta|^2-|a_j+\beta|^2=0,
\]
that is,
\[
 2\Re\bigl(\beta \overline{(a_i-a_j)}\bigr)+|a_i|^2-|a_j|^2=0.
\]
For fixed $i\neq j$, this is a nontrivial affine real linear equation in $\beta\in\mathbb{C}\cong\mathbb{R}^2$, so it defines an affine
real line. Thus the bad values of $\beta$ lie in a finite union of lines. Fix the previously chosen $\lambda\notin \Sigma_1\cup\Sigma_2$.
By Step~1 there exists $\beta_1>0$ such that every $\beta$ with $|\beta|<\beta_1$ preserves the zero configuration in
the disks $D_j$, and by Step~4 together with the nonvanishing verification above,
every $\beta$ with $|\beta|<\beta_*$ keeps all critical values in $K$ above $2\alpha$. Hence,
with
\[
\beta_0:=\min\{\beta_1,\beta_*\}>0,
\]
we may choose $\beta$ with $0<|\beta|<\beta_0$ outside that finite union of lines. Then the numbers
\[
|a_1+\beta|,\dots,|a_N+\beta|
\]
are pairwise distinct. Since the map $t\mapsto -\log t$ is strictly decreasing on
$(0,\infty)$, it follows that the critical values of $u_g$ arising from critical points in
$K$ and larger than $2\alpha$ are pairwise distinct.

\indent This proves the lemma. \end{proof}

\begin{remarkplain}
Lemma~\ref{lem:explicit-perturbation} is only a \emph{local} genericity statement on the fixed compact set $K$. It does not claim anything about critical points of $u_g$ outside $K$. This is sufficient for the argument below, because Lemmas~\ref{lem:component-stability} and~\ref{lem:collar-integral-continuity} show that the relevant perturbed component $V_g$ lies in $K$, so every critical point relevant to Proposition~\ref{prop:generic-spanning-tree} is contained in $K$.
\end{remarkplain}



\section{Stability of the collar integral and of the relevant component}

We retain the notation of the previous section.

\medskip
\begin{lemma}[Component stability]\label{lem:component-stability}
\emph{Suppose $g\to f$ uniformly on $K$ and that, for all sufficiently
close $g$,}
\begin{enumerate}
\item[(a)] \emph{each $D_j$ contains exactly one zero of $g$;}
\item[(b)] \emph{$g$ has no zeros on $K\setminus \bigcup_j D_j$.}
\end{enumerate}
\emph{Then, for all sufficiently small perturbations, there exists a unique connected component $V_g$ of
\[
\{z\in K:-\log|g(z)|>\alpha\}
\]
such that
\[
L_0\subset V_g\subset \operatorname{int}(K)\subset U,
\]
and $V_g$ contains exactly the $m$ zeros of $g$ lying in the disks $D_j$.}
\end{lemma}

\begin{proof} On $L_0$ one has $u\ge 3\alpha/2$, while on $\partial K$ one has $u=\alpha/2$. Uniform convergence implies
\[
 -\log|g|>\alpha \text{ on }L_0,\qquad -\log|g|<\alpha \text{ on }\partial K
\]
for all sufficiently small perturbations. Thus the connected set $L_0$ lies in a unique connected component
\[
 V_g\subset \operatorname{int}(K)
\]
of the open set
\[
\{-\log|g|>\alpha\}\cap \operatorname{int}(K).
\]
Since $-\log|g|<\alpha$ on $\partial K$, this is equivalently a connected component of
$\{-\log|g|>\alpha\}\cap K$.

\indent Fix $j$. Since $\partial D_j\subset L_0\subset V_g$, we have $|g|<e^{-\alpha}$ on $\partial D_j$. By the maximum modulus principle,
\[
 |g|<e^{-\alpha} \text{ throughout } D_j.
\]
Hence $D_j\subset V_g$, so the zero of $g$ inside $D_j$ belongs to $V_g$. There are no other zeros of $g$ in $K$, hence
$V_g$ contains exactly these $m$ zeros.

\indent We next prove uniqueness. Suppose, for contradiction, that $W$ is another connected
component of
\[
\{-\log|g|>\alpha\}\cap K.
\]
Because $-\log|g|<\alpha$ on $\partial K$, every connected component of
$\{-\log|g|>\alpha\}\cap K$ is contained in $\operatorname{int}(K)$; equivalently, $W$ is a
connected component of the open set
\[
\{-\log|g|>\alpha\}\cap \operatorname{int}(K).
\]
Since $V_g$ contains every disk $D_j$, the set $W$ is disjoint from each $D_j$.

We claim that in fact
\[
\overline{W}\cap \partial D_j=\varnothing
\qquad
(1\le j\le m).
\]
Assume otherwise, and choose
\[
p\in \overline{W}\cap \partial D_j
\]
for some $j$. Because $\partial D_j\subset L_0\subset V_g$, and because
$L_0\subset \operatorname{int}(K)$ while $-\log|g|>\alpha$ on $L_0$, there exists $r>0$
such that
\[
B(p,r)\subset \{-\log|g|>\alpha\}\cap \operatorname{int}(K).
\]
Since $p\in \partial D_j$ and $D_j\subset V_g$, the ball $B(p,r)$ meets $V_g$.
Because $p\in \overline{W}$, it also meets $W$. This is impossible, since a connected
subset of the open set $\{-\log|g|>\alpha\}\cap \operatorname{int}(K)$ cannot meet two distinct connected
components of that set. This proves the claim.

Therefore
\[
\overline{W}\subset K\setminus \bigcup_{j=1}^m D_j.
\]
By assumption, $g$ has no zeros on
\[
K \setminus \bigcup_{j=1}^m D_j.
\]
Since $W$ is disjoint from each $D_j$ and $W \Subset \operatorname{int}(K)$, the compact
set $\overline W$ is contained in the open zero-free set
\[
\operatorname{int}(K)\setminus \bigcup_{j=1}^m D_j.
\]
Hence there exists an open neighborhood $N_W$ of $\overline W$ on which $g\neq 0$.
Therefore
\[
-\log|g|
\]
is harmonic on $N_W$, and in particular on a neighborhood of $\overline W$.

We claim that
\[
\partial W \subset \bigl(K\cap\{-\log|g|=\alpha\}\bigr)\cup \partial K.
\]
Indeed, let $p\in \partial W\cap \operatorname{int}(K)$. Since $W$ is a connected
component of the open set $\{-\log|g|>\alpha\}\cap \operatorname{int}(K)$, the point $p$ cannot satisfy
$-\log|g(p)|>\alpha$; otherwise, by continuity, some ball $B(p,r)\subset \operatorname{int}(K)$
would be contained in $\{-\log|g|>\alpha\}\cap \operatorname{int}(K)$ and would meet $W$, forcing
$B(p,r)\subset W$, contrary to $p\in\partial W$. On the other hand, $p$ is a limit point
of $W\subset\{-\log|g|>\alpha\}$, so continuity gives $-\log|g(p)|\ge \alpha$. Hence
$-\log|g(p)|=\alpha$.

If instead $p\in \partial W\cap \partial K$, then the first paragraph of the proof gives
$-\log|g(p)|<\alpha$. Therefore
\[
-\log|g|\le \alpha \qquad \text{on }\partial W.
\]

Since $W$ is bounded and $-\log|g|$ is harmonic on a neighborhood of $\overline W$,
the maximum principle yields
\[
-\log|g|\le \alpha \qquad \text{throughout }W,
\]
contradicting
\[
W\subset \{-\log|g|>\alpha\}.
\]
Thus $V_g$ is the unique connected component of
\[
\{-\log|g|>\alpha\}\cap K.
\]

Finally, because $-\log|g|<\alpha$ on $\partial K$, a path in $\{-\log|g|>\alpha\}$
starting in $V_g$ cannot cross $\partial K$. Hence $V_g$ is also a connected component
of the full superlevel set
\[
\{-\log|g|>\alpha\}.
\]
\end{proof}

\medskip
\begin{lemma}[Continuity of the collar integral]\label{lem:collar-integral-continuity}
\emph{Assume $g\to f$ in $C^1(K)$, and assume that, for all sufficiently close $g$,}
\begin{enumerate}
\item[(a)] \emph{each $D_j$ contains exactly one zero of $g$;}
\item[(b)] \emph{$g$ has no zeros on $K\setminus \bigcup_{j=1}^m D_j$.}
\end{enumerate}
\emph{Let $V_g$ be the corresponding component from Lemma~\ref{lem:component-stability}, and define}
\[
q_g:=\frac1{2\pi}\int_\alpha^{2\alpha} P_g(t)\,dt,
\]
\emph{where}
\[
P_g(t):=\mathcal H^1(\{z\in V_g:u_g(z)=t\})
\qquad\text{for regular }t,
\]
\emph{and $u_g:=-\log|g|$. Then}
\[
q_g\to q
\qquad\text{as }g\to f\text{ in }C^1(K).
\]
\emph{In particular, for all sufficiently small perturbations,}
\[
q_g>\frac q2.
\]
\end{lemma}

\begin{proof}
Because the disks $D_j$ were chosen inside $\{u>5\alpha/2\}$, the compact band
\[
B:=\{z\in U:\alpha/2\le u(z)\le 5\alpha/2\}
\]
is disjoint from $\bigcup_{j=1}^m D_j$. Since all zeros of $g$ in $K$ lie in the disks
$D_j$, the function $u_g=-\log|g|$ is well defined on $B$ for all sufficiently small
perturbations.

Because $f$ has no zeros on
\[
C:=K\setminus \bigcup_{j=1}^m \operatorname{int}(D_j),
\]
there exists $m_C>0$ such that
\[
|f|\ge m_C \qquad\text{on }C.
\]
Since $D_j\subset\{u>5\alpha/2\}$ for each $j$, we also have
\[
M_D:=\max_{z\in \bigcup_{j=1}^m D_j}|f(z)|\le e^{-5\alpha/2}<e^{-2\alpha}.
\]
Choose the perturbation so small that
\[
\|g-f\|_{L^\infty(K)}<\min\{m_C/2,\,e^{-2\alpha}-M_D\}.
\]
Then $g$ has no zeros on $C$, and in fact
\[
|g|\ge m_C/2 \qquad\text{on }C,
\]
while on the union of the root disks we have
\[
|g|\le M_D+\|g-f\|_{L^\infty(K)}<e^{-2\alpha}.
\]
Equivalently,
\[
u_g>2\alpha \qquad\text{on }\bigcup_{j=1}^m D_j.
\]
Hence the entire collar region $\{\alpha<u_g<2\alpha\}\cap K$ is disjoint from the disks $D_j$.

On the zero-free compact set $C$, the formulas
\[u=-\log|f|,\qquad u_g=-\log|g|,
\]
\[
\nabla u=-\nabla(\log|f|),\qquad \nabla u_g=-\nabla(\log|g|)
\]
and the lower bound $|g|\ge m_C/2$ show that $g\to f$ in $C^1(K)$ implies
\[u_g\to u \quad\text{uniformly on }C,
\qquad
\nabla u_g\to \nabla u \quad\text{uniformly on }C.
\]
In particular, since $B\subset C$, we have uniform convergence on $B$ as well. Since $u$
has no critical points on $B$, there exists $c_0>0$ such that
\[
|\nabla u|\ge c_0 \qquad\text{on }B.
\]
Hence, for all sufficiently small perturbations,
\[
|\nabla u_g|\ge c_0/2 \qquad\text{on }B,
\]
so $u_g$ has no critical points on $B$.

We next identify the relevant collar region for $u_g$.

\medskip
\noindent
\emph{Claim.}
\[
\{z\in V_g : \alpha<u_g(z)<2\alpha\}
=
\{z\in K : \alpha<u_g(z)<2\alpha\}
\subset B.
\]
Likewise,
\[
\{z\in U : \alpha<u(z)<2\alpha\}
=
\{z\in K : \alpha<u(z)<2\alpha\}
\subset B.
\]

\smallskip
\noindent
\emph{Proof of claim.}
Let $z\in K$ satisfy
\[
\alpha<u_g(z)<2\alpha.
\]
Since $u_g>2\alpha$ on each disk $D_j$, the point $z$ cannot belong to
\[
\bigcup_{j=1}^m D_j.
\]
Hence $z\in C$. Because $u_g\to u$ uniformly on $C$, for all sufficiently small
perturbations we have
\[
\|u_g-u\|_{L^\infty(C)}<\frac{\alpha}{2}.
\]
Therefore
\[
\frac{\alpha}{2}<u(z)<\frac{5\alpha}{2},
\]
so $z\in B$.

Moreover, $u_g(z)>\alpha$, and Lemma~\ref{lem:component-stability} says that $V_g$ is the
unique connected component of
\[
\{u_g>\alpha\}\cap K.
\]
Thus $z\in V_g$. We have proved that
\[
\{z\in K : \alpha<u_g(z)<2\alpha\}\subset V_g\cap B.
\]
The reverse inclusion is immediate, so
\[
\{z\in V_g : \alpha<u_g(z)<2\alpha\}
=
\{z\in K : \alpha<u_g(z)<2\alpha\}
\subset B.
\]

For the unperturbed function $u$, recall that
\[
K=\{u\ge \alpha/2\}.
\]
Hence any point $z\in U$ with $\alpha<u(z)<2\alpha$ automatically lies in $K$, and
therefore in $B$. Thus
\[
\{z\in U : \alpha<u(z)<2\alpha\}
=
\{z\in K : \alpha<u(z)<2\alpha\}
\subset B.
\]
This proves the claim.
\medskip

Applying the coarea formula on the band $B$, we obtain
\[
2\pi q=\int_B 1_{\{\alpha<u<2\alpha\}}\,|\nabla u|\,dA,
\]
and
\[
2\pi q_g=\int_B 1_{\{\alpha<u_g<2\alpha\}}\,|\nabla u_g|\,dA.
\]

Set
\[
E:=\{z\in B:\alpha<u(z)<2\alpha\},
\qquad
E_g:=\{z\in B:\alpha<u_g(z)<2\alpha\}.
\]
Because $u_g\to u$ uniformly on $B$, we have
\[
1_{E_g}(z)\to 1_E(z)
\]
for every $z\in B$ such that $u(z)\notin\{\alpha,2\alpha\}$. The exceptional set is
contained in
\[
\{u=\alpha\}\cup\{u=2\alpha\},
\]
which has planar measure zero by Lemma~\ref{lem:level-set-measure-zero}. Thus $1_{E_g}\to 1_E$ almost everywhere
on $B$.

Also, for $g$ sufficiently close to $f$, uniform convergence of the gradients gives a
constant $M_B<\infty$ such that
\[
|\nabla u_g|\le M_B \qquad\text{on }B.
\]
Therefore the integrands
\[
1_{E_g}|\nabla u_g|
\]
are dominated by the integrable function $M_B1_B$, and they converge almost
everywhere on $B$ to
\[
1_E|\nabla u|.
\]
By dominated convergence,
\[
q_g\to q.
\]
In particular, $q_g>q/2$ for all sufficiently small perturbations.
\end{proof}


\section{Proof of the main theorem}

We now complete the argument.

\medskip
\begin{lemma}[A tree lemma]\label{lem:tree-lemma}
\emph{Let $T$ be a finite tree of total length $L$, and let $x_1,\dots,x_M$ be
pairwise distinct marked points of $T$, with $M\ge 2$. Then there exist
distinct indices $i\neq j$ such that the unique subpath joining $x_i$ to $x_j$
has length at most $2L/M$.}
\end{lemma}

\begin{proof}
Subdivide edges at the marked points; this does not change the total length and
turns the marked points into marked vertices. Let $T'\subseteq T$ be the minimal
subtree spanning the marked vertices, and let $L'=\len(T')$. Then $L'\le L$,
and every leaf of $T'$ is marked.

Doubling every edge of $T'$ produces an Eulerian graph of total length $2L'$.
Choose an Euler tour that starts at a marked vertex, and let
\[
y_1,\dots,y_M
\]
be the marked vertices listed in the order of their first appearance along the
tour. Let
\[
0=t_1<t_2<\cdots<t_M<2L'
\]
be their first-appearance times, and set
\[
t_{M+1}:=2L'.
\]
Then
\[
\sum_{r=1}^M (t_{r+1}-t_r)=2L'.
\]

For each $r=1,\dots,M-1$, the portion of the Euler tour between times $t_r$ and
$t_{r+1}$ is a walk in $T'$ joining $y_r$ to $y_{r+1}$, so
\[
d_{T'}(y_r,y_{r+1})\le t_{r+1}-t_r.
\]
Likewise, the final portion of the tour from time $t_M$ to time $2L'$ joins
$y_M$ back to $y_1$, so
\[
d_{T'}(y_M,y_1)\le 2L'-t_M=t_{M+1}-t_M.
\]
Hence
\[
\sum_{r=1}^{M-1} d_{T'}(y_r,y_{r+1}) + d_{T'}(y_M,y_1)
\le 2L' \le 2L.
\]
Therefore at least one of these $M$ distances is at most $2L/M$, which proves
the lemma.
\end{proof}

\begin{proof}[Proof of Theorem~\ref{thm:main}]
If $n=1$, then the unique root may be joined to itself by the constant curve, so the conclusion is immediate. If $f$ has a multiple zero, the conclusion is also immediate. Thus we may assume that $f$ is squarefree and $n\ge 2$.

\indent By the component lemma, there exists a connected component $U\subset \Lambda_f$ containing exactly
$m\ge 2$ zeros of $f$. Let $u=-\log|f|$ on $U$, let $\alpha$ and $q$ be as in Section~3, and choose the root disks
$D_1,\dots,D_m$ as in Section~5, so small that
\begin{equation}
\sum_{j=1}^m \operatorname{diam}(D_j)<\frac{q}{4m}.
\tag{17}
\end{equation}
Choose $(\lambda,\beta)$ as in Lemma~\ref{lem:explicit-perturbation} so small that
\[
\|g_{\lambda,\beta}-f\|_{C^1(K)}<\varepsilon_0,
\]
where $\varepsilon_0>0$ is small enough for Lemmas~\ref{lem:component-stability} and~\ref{lem:collar-integral-continuity} to apply, and set
\[
g:=g_{\lambda,\beta}.
\]
By Lemma~\ref{lem:explicit-perturbation}(1), each disk $D_j$ contains exactly one simple zero $b_j$ of $g$, and
$g$ has no zeros on
\[
K\setminus \bigcup_{j=1}^m D_j.
\]
Because the disks $D_j$ lie in $\{u>5\alpha/2\}$, they are disjoint from the fixed band
\[
B=\{z\in U:\alpha/2\le u\le 5\alpha/2\},
\]
so $g$ has no zeros on $B$ as well. Let $V_g$ be the corresponding component from
Lemma~\ref{lem:component-stability}. Then $V_g$ is a connected component of the full superlevel set
\[
\{-\log|g|>\alpha\},
\]
it satisfies $V_g\subset \operatorname{int}(K)\subset U$, it contains exactly the $m$
perturbed roots $b_1,\dots,b_m$ lying in the disks $D_j$, and by Lemma~\ref{lem:collar-integral-continuity} we have
\[
q_g>q/2.
\]

Applying Corollary~\ref{cor:component-integral-bound} with
\[
h=g,
\qquad
v=u_g:=-\log|g|,
\qquad
c=\alpha,
\qquad
W=V_g,
\]
gives
\[
\frac1{2\pi}\int_\alpha^\infty P_g(t)\,dt\le m.
\]
Hence
\[
\frac1{2\pi}\int_{2\alpha}^\infty P_g(t)\,dt
=
\frac1{2\pi}\int_{\alpha}^\infty P_g(t)\,dt
-
\frac1{2\pi}\int_{\alpha}^{2\alpha} P_g(t)\,dt
\le m-q_g
<
m-\frac q2.
\]

Since all zeros of $g$ in $V_g$ are the simple roots $b_1,\dots,b_m$, Proposition~\ref{prop:generic-spanning-tree} applies to $u_g=-\log|g|$ on $V_g$ with base level $c=\alpha$. By Lemma~\ref{lem:component-stability} we have $V_g\subset K$, so every critical point relevant to Proposition~\ref{prop:generic-spanning-tree} lies in $K$. Therefore Lemma~\ref{lem:explicit-perturbation}(2) gives the absence of critical values in $[\alpha,2\alpha]$, while Lemma~\ref{lem:explicit-perturbation}(3)--(4) give the Morse and distinct-critical-value hypotheses above $2\alpha$. Choose
\[
\varepsilon:=\frac q8.
\]
Then Proposition~\ref{prop:generic-spanning-tree} yields a finite tree \(G\subset V_g\subset U\) containing the \(m\) perturbed roots
and satisfying
\[
\len(G)\le \frac1{2\pi}\int_{2\alpha}^\infty P_g(t)\,dt+\frac q8
<
m-\frac q2+\frac q8
=
m-\frac{3q}{8}.
\]

By Lemma~\ref{lem:tree-lemma}, since the $m$ perturbed roots are pairwise
distinct and $m\ge 2$, there exist distinct indices $r\neq s$ such that
the points \(b_r\in D_r\) and \(b_s\in D_s\) are joined in \(G\) by a subpath of
length at most
\[
\frac{2}{m}\len(G)
<
\frac{2}{m}\left(m-\frac{3q}{8}\right)
=
2-\frac{3q}{4m}.
\tag{18}
\]

Connect \(b_r\) to the original zero \(a_r\in D_r\) and \(b_s\) to the original zero \(a_s\in D_s\) by
straight line segments. Because each $D_j$ is an ordinary Euclidean disk contained in $U$,
these straight segments lie in
\[
D_r\cup D_s \subset U \subset \Lambda_f.
\]
The total additional length is at most
\[
\operatorname{diam}(D_r)+\operatorname{diam}(D_s)
\le
\sum_{j=1}^m \operatorname{diam}(D_j)
<
\frac{q}{4m}
\]
by \((17)\). Let \(\Gamma\) be the union of the subpath of \(G\) joining \(b_r\) to \(b_s\) and the two straight
segments \([a_r,b_r]\) and \([b_s,a_s]\). Then \(\Gamma\) is a connected finite graph contained in
\[
U\subset \Lambda_f.
\]
Hence \(\Gamma\) contains a simple arc \(\gamma\) joining \(a_r\) to \(a_s\), and
\[
\len(\gamma)\le \len(\Gamma).
\]
Therefore
\[
\len(\gamma)
<
\left(2-\frac{3q}{4m}\right)+\frac{q}{4m}
=
2-\frac{q}{2m}
<2.
\]
Thus \(\gamma\subset \Lambda_f\) is a curve joining two zeros of \(f\) and having length \(<2\).
This proves the theorem.
\end{proof}

\begin{thebibliography}{9}
\bibitem{Crane2004}
E.~Crane, \emph{The area of polynomial images and preimages}, Bull. Lond. Math. Soc. \textbf{36} (2004),
786--792.

\bibitem{EHP58}
P.~Erd\H{o}s, F.~Herzog, and G.~Piranian, \emph{Metric properties of polynomials}, J. Analyse Math. \textbf{6}
(1958), 123--148.

\bibitem{GR2024}
S.~Ghosh and K.~Ramachandran, \emph{Number of components of polynomial lemniscates: a problem
of Erd\H{o}s, Herzog, and Piranian}, J. Math. Anal. Appl. \textbf{540} (2024), no.~1, 128571.

\bibitem{Milnor1963}
J.~Milnor, \emph{Morse Theory}, Annals of Mathematics Studies, No.~51, Princeton University Press,
Princeton, NJ, 1963.
\end{thebibliography}

\end{document}
